% =============================================================================
%  main.tex -- składa raport: strona tytułowa, streszczenie, rozdziały, dodatki
%
%  Budowanie (z tego katalogu):   latexmk -pdf main.tex
%  lub ręcznie:                   pdflatex main && bibtex main && pdflatex main && pdflatex main
%
%  Tabele z danymi w tables/ są generowane z wyników przebiegu:
%      python Report/tools/make_tables.py results/haiku-v1
% =============================================================================
\documentclass[11pt,a4paper,oneside]{report}

\newcommand{\ReportTitle}{Czy paradygmat programowania ma znaczenie?\\[2pt]
  \Large Ocena kodu generowanego przez LLM w językach imperatywnych,\\
  obiektowych, funkcyjnych i logicznych}
\newcommand{\ShortTitle}{Czy paradygmat programowania ma znaczenie?}

% Autorzy. TODO(autorzy): uzupełnić imiona i nazwiska pozostałych członków zespołu
% (strona tytułowa i metadane PDF poniżej korzystają z tych makr).
\newcommand{\AuthorOne}{Hubert Wajda}
\newcommand{\AuthorTwo}{Stanisław Lemański}
\newcommand{\AuthorThree}{Bartosz Kopania}
\newcommand{\AuthorFour}{[Imię i nazwisko 4]}
\newcommand{\AuthorFive}{[Imię i nazwisko 5]}
\newcommand{\AuthorSix}{[Imię i nazwisko 6]}

% =============================================================================
%  preamble.tex -- packages, colours, code styles and macros of the report
% =============================================================================

\PassOptionsToPackage{table}{xcolor}

% ---------------------------------------------------------------- fonts & text
\usepackage[T1]{fontenc}
\usepackage[utf8]{inputenc}
\usepackage[english,polish]{babel}      % Polish is the main language
\usepackage{icomma}                       % decimal comma in maths: $0,05$
\usepackage{newtxtext,newtxmath}          % Times-like text and maths
\usepackage[varqu,varl,scaled=0.94]{zi4}  % Inconsolata for code
\usepackage[a4paper,margin=2.5cm,headheight=14pt]{geometry}
\usepackage{microtype}
\usepackage{setspace}
\setstretch{1.08}
\usepackage{pifont}                       % check / cross marks
\usepackage{csquotes}

% ---------------------------------------------------------------- floats & tables
\usepackage{xcolor}
\usepackage{graphicx}
\usepackage{booktabs,tabularx,multirow,makecell,longtable,array}
\usepackage{threeparttable}
\usepackage[detect-all]{siunitx}
\usepackage{enumitem}
\usepackage[font=small,labelfont=bf,skip=6pt]{caption}
\usepackage{subcaption}
\usepackage{float}
\setlist{itemsep=2pt,topsep=4pt}
\renewcommand{\arraystretch}{1.12}
\newcolumntype{L}[1]{>{\raggedright\arraybackslash}p{#1}}
\newcolumntype{C}[1]{>{\centering\arraybackslash}p{#1}}
\newcolumntype{Y}{>{\raggedright\arraybackslash}X}

% ---------------------------------------------------------------- figures
\usepackage{tikz}
\usetikzlibrary{arrows.meta,positioning,fit,calc,shapes.geometric,backgrounds}
\usepackage{pgfplots}
\pgfplotsset{compat=1.18}

% ---------------------------------------------------------------- headings & page style
\usepackage{titlesec}
\titleformat{\chapter}[hang]{\normalfont\LARGE\bfseries}{\thechapter}{0.7em}{}
\titlespacing*{\chapter}{0pt}{-12pt}{16pt}
\titleformat{\section}{\normalfont\Large\bfseries}{\thesection}{0.7em}{}
\titleformat{\subsection}{\normalfont\large\bfseries}{\thesubsection}{0.7em}{}
\titleformat{\paragraph}[runin]{\normalfont\normalsize\bfseries}{}{0pt}{}[.]
\titlespacing*{\paragraph}{0pt}{6pt plus 2pt}{0.6em}
\setcounter{secnumdepth}{2}
\setcounter{tocdepth}{1}

\usepackage{fancyhdr}
\pagestyle{fancy}
\fancyhf{}
\fancyhead[L]{\small\itshape\ShortTitle}
\fancyhead[R]{\small\itshape\nouppercase{\leftmark}}
\fancyfoot[C]{\small\thepage}
\renewcommand{\headrulewidth}{0.4pt}
\renewcommand{\chaptermark}[1]{\markboth{\thechapter.\ #1}{}}
\fancypagestyle{plain}{\fancyhf{}\fancyfoot[C]{\small\thepage}\renewcommand{\headrulewidth}{0pt}}

% ---------------------------------------------------------------- colours
% Chart palette: categorical slots 1-4 of the validated reference palette (adjacent-pair
% CVD-safe in this order) and the single-hue blue ramp for ordinal/sequential encodings.
\definecolor{ink}{HTML}{0B0B0B}
\definecolor{inksec}{HTML}{52514E}
\definecolor{inkmuted}{HTML}{8A8984}
\definecolor{gridline}{HTML}{E4E3DF}
\definecolor{neutral}{HTML}{C3C2B7}
\definecolor{neutralbg}{HTML}{F0EFEC}
\definecolor{cat1}{HTML}{2A78D6}   % blue    -- imperative
\definecolor{cat2}{HTML}{EB6834}   % orange  -- object-oriented
\definecolor{cat3}{HTML}{1BAF7A}   % aqua    -- functional
\definecolor{cat4}{HTML}{EDA100}   % yellow  -- logic
\definecolor{seq050}{HTML}{E6F0FC}
\definecolor{seq150}{HTML}{B7D3F6}
\definecolor{seq250}{HTML}{86B6EF}
\definecolor{seq300}{HTML}{6DA7EC}
\definecolor{seq400}{HTML}{3987E5}
\definecolor{seq500}{HTML}{256ABF}
\definecolor{seq550}{HTML}{1C5CAB}
\definecolor{seq600}{HTML}{184F95}
\definecolor{seq650}{HTML}{104281}
% code highlighting
\definecolor{codebg}{HTML}{F7F7F5}
\definecolor{coderule}{HTML}{C9C8C2}
\definecolor{codekw}{HTML}{1C5CAB}
\definecolor{codecm}{HTML}{5F7A5F}
\definecolor{codestr}{HTML}{9A3E12}
\definecolor{todo}{HTML}{B3261E}

% ---------------------------------------------------------------- code listings
\usepackage{listings}
\lstset{
  basicstyle=\ttfamily\footnotesize,
  columns=fullflexible, keepspaces=true,
  breaklines=true, breakatwhitespace=false,
  postbreak=\mbox{\textcolor{inkmuted}{$\hookrightarrow$}\space},
  frame=lines, framesep=4pt, rulecolor=\color{coderule},
  backgroundcolor=\color{codebg},
  numbers=left, numberstyle=\tiny\color{inkmuted}, numbersep=6pt, xleftmargin=1.6em,
  keywordstyle=\color{codekw}\bfseries,
  commentstyle=\color{codecm},          % Inconsolata has no italic shape
  stringstyle=\color{codestr},
  showstringspaces=false, upquote=true, tabsize=4,
  captionpos=t, aboveskip=8pt, belowskip=6pt,
  extendedchars=true, inputencoding=utf8,
  literate={≤}{{$\le$}}1 {≥}{{$\ge$}}1 {—}{{---}}1 {–}{{--}}1 {→}{{$\rightarrow$}}1
           {×}{{$\times$}}1 {…}{{\ldots}}1,
}
\lstdefinestyle{prompt}{
  language={}, numbers=none, xleftmargin=0pt,
  basicstyle=\ttfamily\footnotesize, backgroundcolor=\color{codebg},
  keywordstyle={}, commentstyle={}, stringstyle={},
}
\lstdefinestyle{plain}{language={}, keywordstyle={}, commentstyle={}, stringstyle={}}
\lstdefinelanguage{Smalltalk}{
  morekeywords={self,super,nil,true,false,thisContext},
  morecomment=[s]{"}{"}, morestring=[b]',
  sensitive=true,
}
\lstdefinelanguage{JSON}{
  morestring=[b]", stringstyle=\color{codestr},
  literate={:}{{\textcolor{inksec}{:}}}1 {,}{{\textcolor{inksec}{,}}}1,
}
\renewcommand{\lstlistingname}{Listing}
\renewcommand{\lstlistlistingname}{Spis listingów}

% ---------------------------------------------------------------- references
\usepackage[numbers,sort&compress]{natbib}
\usepackage[pdfusetitle,bookmarksnumbered]{hyperref}
\hypersetup{colorlinks=true,linkcolor=seq600,citecolor=seq600,urlcolor=seq600}
\usepackage[nameinlink]{cleveref}
% cleveref has no Polish definitions; abbreviated forms avoid Polish declension in running text.
\crefname{chapter}{rozdz.}{rozdz.}             \Crefname{chapter}{Rozdz.}{Rozdz.}
\crefname{appendix}{dod.}{dod.}                \Crefname{appendix}{Dod.}{Dod.}
\crefname{section}{podrozdz.}{podrozdz.}       \Crefname{section}{Podrozdz.}{Podrozdz.}
\crefname{subsection}{podrozdz.}{podrozdz.}    \Crefname{subsection}{Podrozdz.}{Podrozdz.}
\crefname{subappendix}{dod.}{dod.}             \Crefname{subappendix}{Dod.}{Dod.}
\crefname{table}{tab.}{tab.}                   \Crefname{table}{Tab.}{Tab.}
\crefname{figure}{rys.}{rys.}                  \Crefname{figure}{Rys.}{Rys.}
\crefname{subfigure}{rys.}{rys.}               \Crefname{subfigure}{Rys.}{Rys.}
\crefname{lstlisting}{list.}{list.}            \Crefname{lstlisting}{List.}{List.}
\crefname{enumi}{pkt}{pkt}                     \Crefname{enumi}{Pkt}{Pkt}
\crefname{rq}{RQ}{RQ}
\newcommand{\crefrangeconjunction}{--}
\newcommand{\crefpairconjunction}{ i }
\newcommand{\crefmiddleconjunction}{, }
\newcommand{\creflastconjunction}{ i }
\newcommand{\crefpairgroupconjunction}{ oraz }
\newcommand{\crefmiddlegroupconjunction}{, }
\newcommand{\creflastgroupconjunction}{ oraz }
\urlstyle{same}

% ---------------------------------------------------------------- macros
% Expandable input for generated table bodies: LaTeX's \input is not expandable and breaks
% when the included rows are followed by a rule (Misplaced \noalign).
\makeatletter
\newcommand*{\tabinput}[1]{\@@input #1.tex }
\makeatother
\newcommand{\code}[1]{\texttt{\detokenize{#1}}}       % inline identifiers (no # inside)
\newcommand{\file}[1]{\texttt{\detokenize{#1}}}
\newcommand{\cmark}{\ding{51}}
\newcommand{\xmark}{\ding{55}}
\newcommand{\haiku}{Claude Haiku~5.5}
\newcommand{\sonnet}{Claude Sonnet~5.5}
\newcommand{\RQ}[1]{\textbf{RQ#1}}
\newcommand{\usd}[1]{\$#1}

% Result cells of the language x task matrix (tables/matrix.tex):
%   \resfirst{final}           solved at the first attempt
%   \resrep{first}{final}{k}   solved in repair round k
%   \resfail{first}{final}     not solved after all repair rounds
\newcommand{\resfirst}[1]{#1\,\textcolor{seq600}{\cmark}}
\newcommand{\resrep}[3]{#1$\to$#2\,\textcolor{seq600}{\cmark}\textsuperscript{#3}}
\newcommand{\resfail}[2]{\cellcolor{neutralbg}#1$\to$#2\,\textcolor{todo}{\xmark}}

% Visible notes for information the authors still have to supply.
% Set \showtodosfalse in main.tex to hide them in a final build.
\newif\ifshowtodos
\showtodostrue
\newcommand{\todo}[1]{\ifshowtodos{\color{todo}\sffamily\small[\textbf{TODO:} #1]}\fi}

% Research-question box
\usepackage[most]{tcolorbox}
\newtcolorbox{rqbox}{colback=seq050,colframe=seq500,boxrule=0.5pt,arc=1pt,
  left=6pt,right=6pt,top=4pt,bottom=4pt,fontupper=\small}
\newtcolorbox{findingbox}[1][]{colback=neutralbg,colframe=inksec,boxrule=0.4pt,arc=1pt,
  left=6pt,right=6pt,top=4pt,bottom=4pt,fontupper=\small,title={#1},
  coltitle=ink,colbacktitle=neutralbg,fonttitle=\bfseries\small,attach title to upper={\ }}

% pgfplots defaults shared by all charts
\pgfplotsset{
  reportchart/.style={
    axis line style={draw=inkmuted},
    tick style={draw=inkmuted},
    every tick label/.append style={font=\footnotesize, text=inksec},
    label style={font=\footnotesize, text=inksec},
    legend style={font=\footnotesize, draw=none, fill=none, text=ink},
    legend cell align=left,
    grid style={draw=gridline, line width=0.4pt},
    major grid style={draw=gridline},
    /pgf/number format/use comma,
    /pgf/number format/1000 sep={\,},
  },
}


\hypersetup{
  pdftitle={Czy paradygmat programowania ma znaczenie? Ocena kodu generowanego przez LLM w różnych paradygmatach},
  pdfauthor={\AuthorOne, \AuthorTwo, \AuthorThree, \AuthorFour, \AuthorFive, \AuthorSix},
  pdfkeywords={duże modele językowe, generowanie kodu, paradygmaty programowania, benchmark, LLM jako sędzia},
}

\begin{document}

% ------------------------------------------------------------------ strona tytułowa
\begin{titlepage}
  \thispagestyle{empty}
  \centering
  \vspace*{0.6cm}
  {\small\scshape Raport z projektu badawczego\par}
  \vspace{1.2cm}
  {\huge\bfseries \ReportTitle\par}
  \vspace{1.2cm}
  {\large \AuthorOne \quad \AuthorTwo \quad \AuthorThree\par}
  \vspace{0.2cm}
  {\large \AuthorFour \quad \AuthorFive \quad \AuthorSix\par}
  \vspace{0.35cm}
  % TODO(autorzy): potwierdzić afiliację i kolejność autorów.
  {\normalsize\itshape Politechnika Gdańska\par}
  \vspace{0.25cm}
  {\small \today\par}
  \vspace{1.1cm}
  \begin{minipage}{0.86\textwidth}
    % Streszczenie -- składane na stronie tytułowej przez main.tex
\begin{center}\textbf{Streszczenie}\end{center}
\vspace{-4pt}
{\small
Duże modele językowe (LLM) ocenia się zwykle na benchmarkach generowania kodu zdominowanych przez
Pythona i kilka popularnych języków, dlatego w dużej mierze nie wiadomo, czy sam \emph{paradygmat
programowania} wpływa na poprawność i idiomatyczność generowanego kodu. Przedstawiamy
\emph{LLM Paradigm Benchmark}: dziesięć zadań algorytmicznych, od wykrywania kolizji za pomocą sumy
Minkowskiego i algorytmu Dijkstry po eertree, dopasowanie Gotoha, heurystykę Christofidesa, B-drzewa
w ujęciu CLRS i selekcję BFPRT, każde z \emph{tymi samymi} 20 ukrytymi testami jednostkowymi
w siedmiu językach z czterech paradygmatów: imperatywnego (C, Pascal), obiektowego (Java, Smalltalk),
funkcyjnego (Haskell, Scala) i logicznego (Prolog). Zautomatyzowany system ewaluacyjny przekazuje
modelowi specyfikację i kontrakt kodu, uruchamia odpowiedź w izolowanych kontenerach Docker (tam,
gdzie to możliwe, pod Valgrindem), zwraca diagnostykę kompilatora i testów w maksymalnie trzech
rundach samonaprawy, a na koniec prosi zaślepionego sędziego LLM o ocenę idiomatyczności według
rubryki właściwej dla danego języka. W pełnym przebiegu modelu \haiku{} (70~zadań, 1\,400~testów),
ocenianym przez \sonnet{}, model rozwiązał 80,0\,\% zadań za pierwszym podejściem i 92,9\,\% po
samonaprawie (87,8\,\% $\to$ 98,2\,\% testów). Różnice w poprawności są większe między językami niż
między paradygmatami i żadna z nich nie jest istotna statystycznie przy tej wielkości próby.
Idiomatyczność wyraźniej rozróżnia paradygmaty: kod funkcyjny uzyskał najwyższą ocenę (3,90 na 5),
natomiast kod obiektowy (3,45) i logiczny (3,40) był często pisany w stylu proceduralnym. Przebieg
ujawnił też dwa defekty w zestawie testów dla Prologu. Raportujemy wreszcie negatywne wyniki badań
pilotażowych: zadania w stylu HumanEval okazały się zbyt łatwe, by rozróżnić paradygmaty, a lokalnie
uruchamiane modele o otwartych wagach były zbyt wolne i zbyt często otaczały kod prozą, by nadawały
się do zautomatyzowanego potoku.\par}

\vspace{6pt}
{\small\noindent\textbf{Słowa kluczowe:} duże modele językowe; generowanie kodu; paradygmaty
programowania; benchmark; samonaprawa; LLM jako sędzia; idiomatyczność.\par}

  \end{minipage}
  \vfill
  {\footnotesize Kod, zadania benchmarku i surowe wyniki: repozytorium \texttt{llm-paradigm-benchmark},
   przebieg \texttt{haiku-v1} (commity \texttt{072ea49} i \texttt{dc7f1d4}).\par}
\end{titlepage}

% ------------------------------------------------------------------ część wstępna
\pagenumbering{roman}
\tableofcontents
\clearpage
\pagenumbering{arabic}

% ------------------------------------------------------------------ rozdziały
% =============================================================================
\chapter{Wprowadzenie}
\label{ch:introduction}
% =============================================================================

Duże modele językowe (LLM) stały się codziennym narzędziem programistów. Uzupełniają kod
w edytorach, odpowiadają na pytania o biblioteki i coraz częściej piszą całe funkcje lub moduły na
podstawie krótkiego opisu. Ocena tych umiejętności podążała tą samą drogą. Wczesne benchmarki, takie
jak HumanEval~\cite{chen2021codex} i MBPP~\cite{austin2021mbpp}, mierzą, jak często model tworzy
funkcję w Pythonie, która przechodzi kilka testów jednostkowych. Późniejsze prace przetłumaczyły te
problemy na inne języki~\cite{cassano2023multiple,zheng2023codegeex,athiwaratkun2023mbxp}, dodały
mocniejsze testy~\cite{liu2023evalplus} lub przeszły do zadań na poziomie całych
repozytoriów~\cite{jimenez2024swebench}.

Niemal wszystkie te prace traktują język programowania jako szczegół powierzchniowy: zadanie
pozostaje takie samo, zmienia się tylko składnia. Języki różnią się jednak nie tylko składnią.
\emph{Paradygmat programowania}~\cite{floyd1979paradigms,vanroy2004ctm}, czy to imperatywny,
obiektowy, funkcyjny czy logiczny, określa, jak rozkłada się rozwiązanie na części, jak wyraża się
stan i przepływ danych, jak reprezentuje się brak wartości i błędy oraz kto zarządza pamięcią.
Algorytm Dijkstry napisany z modyfikowalnym kopcem binarnym w C, z niemodyfikowalnymi mapami
i zwinięciami (ang.~\emph{folds}) w Haskellu albo z relacjami i nawracaniem w Prologu to trzy różne
ćwiczenia. Dane treningowe modeli LLM są ponadto silnie przesunięte w stronę kilku popularnych
języków~\cite{kocetkov2022stack}, a modele radzą sobie wyraźnie gorzej z językami o małych zasobach
danych~\cite{cassano2023multiple,cassano2024knowledge}. Smalltalk, Pascal i Prolog są takimi
językami pod każdym względem.

W niniejszym raporcie pytamy, czy te różnice mają znaczenie dla kodu generowanego przez LLM.
Utrzymujemy stałe \emph{zadanie} i \emph{testy}, a zmieniamy tylko język, a wraz z nim paradygmat.
Następnie mierzymy poprawność funkcjonalną, zdolność modelu do naprawy własnych błędów na podstawie
informacji zwrotnej z kompilatora i testów oraz idiomatyczność kodu końcowego. Idiomatyczność jest
istotna, ponieważ rozwiązanie może przejść wszystkie testy, a mimo to być „C napisanym
w Haskellu”, czyli kodem, który programista utrzymujący kod w tym języku by odrzucił.

\section{Pytania badawcze}
\label{sec:rqs}

\begin{rqbox}
\begin{description}[leftmargin=3.2em,labelwidth=2.8em,itemsep=3pt]
  \item[\RQ{1}] \textbf{Poprawność.} Jak różni się poprawność funkcjonalna kodu generowanego przez
    LLM za pierwszym podejściem między językami i paradygmatami, gdy zadanie i testy pozostają
    stałe?
  \item[\RQ{2}] \textbf{Samonaprawa.} W jakim stopniu luka ta się zmniejsza, gdy model otrzymuje
    komunikaty kompilatora, logi testów i nazwy niezaliczonych testów, i czy ta zdolność zależy od
    paradygmatu?
  \item[\RQ{3}] \textbf{Idiomatyczność.} Czy wygenerowany kod jest idiomatyczny dla docelowego
    języka i paradygmatu w ocenie zaślepionego sędziego LLM posługującego się stałą rubryką?
  \item[\RQ{4}] \textbf{Praktyczność.} Ile kosztuje taka ewaluacja i jakie przeszkody praktyczne
    (wykonywanie instrukcji, format odpowiedzi, sposób hostowania modelu) wpływają na jej
    trafność?
\end{description}
\end{rqbox}

\section{Wkład}

\begin{enumerate}
  \item \textbf{Benchmark obejmujący wiele paradygmatów.} Dziesięć zadań algorytmicznych, każde
    z dwoma algorytmami i 20 testami jednostkowymi, zaimplementowanych w siedmiu językach
    reprezentujących cztery paradygmaty (łącznie 1\,400 przypadków testowych). Każde zadanie zawiera
    specyfikację, kontrakt, szkielet (ang.~\emph{stub}) i ukryty zestaw testów. Każdy język ma
    własny obraz Docker, a błędy pamięci są sprawdzane Valgrindem tam, gdzie język na to pozwala
    (\cref{ch:benchmark}).
  \item \textbf{System ewaluacyjny.} Program w Pythonie, który generuje prompty z szablonów dla
    poszczególnych języków, wywołuje model, ocenia odpowiedź w kontenerach, klasyfikuje wynik do
    jednej z dziewięciu kategorii, prowadzi pętlę samonaprawy, ocenia idiomatyczność za pomocą
    zaślepionego sędziego LLM i agreguje raport. Przebiegi można wznawiać, są w pełni śledzone,
    a ich koszt jest kontrolowany (\cref{ch:harness}).
  \item \textbf{Badanie empiryczne.} Pełny przebieg modelu \haiku{} na wszystkich 70 parach
    język--zadanie, oceniony przez \sonnet{}, wraz z analizą poprawności, rodzajów błędów,
    samonaprawy, idiomatyczności i kosztów (\cref{ch:results,ch:discussion}).
  \item \textbf{Wyniki negatywne i wnioski.} Zadania w stylu HumanEval okazały się zbyt łatwe, by
    rozróżnić paradygmaty, a lokalne modele o otwartych wagach były niepraktyczne ze względu na
    szybkość i format odpowiedzi (\cref{ch:preliminary}). Przebieg ujawnił też dwa defekty
    w zestawie testów dla Prologu (\cref{sec:benchmark-defects}).
\end{enumerate}

\section{Struktura raportu}

\Cref{ch:background} przedstawia paradygmaty oraz prace pokrewne dotyczące benchmarków generowania
kodu, samonaprawy i sędziów LLM. \Cref{ch:preliminary} opisuje nasze eksperymenty pilotażowe
i wyciągnięte z nich wnioski. \Cref{ch:benchmark} prezentuje benchmark, a \cref{ch:harness} system
ewaluacyjny, w tym konstrukcję promptów. \Cref{ch:results} zawiera wyniki, \cref{ch:discussion}
ich omówienie wraz z zagrożeniami dla trafności, a \cref{ch:conclusions} podsumowanie. Dodatki
zawierają pełne prompty, rubrykę sędziego i wyniki dla poszczególnych zadań.

% =============================================================================
\chapter{Tło i prace pokrewne}
\label{ch:background}
% =============================================================================

\section{Paradygmaty programowania}
\label{sec:paradigms}

Paradygmat programowania to sposób strukturyzowania obliczeń, który język wspiera i do którego
zachęca~\cite{floyd1979paradigms,vanroy2004ctm}. Rozważamy cztery klasyczne paradygmaty i tam,
gdzie to możliwe, wybieramy po dwa języki na paradygmat. Jeden z nich jest szeroko stosowany,
a drugi jest mniej popularnym, „czystszym” przedstawicielem. Pozwala to oddzielić wpływ paradygmatu
od wpływu popularności pojedynczego języka.

\paragraph{Imperatywny (proceduralny): C i Pascal} Programy są ciągami instrukcji modyfikujących
stan. Dane to rekordy i tablice, a przepływ sterowania opiera się na pętlach i instrukcjach
warunkowych. W C pamięcią zarządza się ręcznie za pomocą \code{malloc}/\code{free}, więc wycieki
i niepoprawne dostępy stanowią odrębną klasę defektów. Pascal wprowadza silniejsze typowanie
i zarządzane tablice dynamiczne, ale zachowuje \code{New}/\code{Dispose} dla wskaźników.

\paragraph{Obiektowy: Java i Smalltalk} Zachowanie znajduje się w obiektach, które hermetyzują stan
i komunikują się za pomocą komunikatów. Java jest typowana statycznie i ma bogatą bibliotekę
standardową. Smalltalk to typowany dynamicznie język, w którym „wszystko jest obiektem”, a nawet
struktury sterujące (\code{ifTrue:}, \code{do:}) są komunikatami wysyłanymi do obiektów i bloków.

\paragraph{Funkcyjny: Haskell i Scala} Obliczenie polega na wartościowaniu wyrażeń bez efektów
ubocznych. Haskell jest czysty i leniwy. Brak wartości i błędy modeluje algebraicznymi typami danych
(\code{Maybe}, \code{Either}) i preferuje zwinięcia, dopasowanie wzorców oraz trwałe struktury
danych. Scala łączy cechy funkcyjne i obiektowe na JVM, oferując niemodyfikowalne kolekcje, klasy
przypadków (ang.~\emph{case classes}) i \code{Option}.

\paragraph{Logiczny: Prolog} Program jest zbiorem relacji (faktów i reguł), a obliczenie polega na
poszukiwaniu dowodu za pomocą unifikacji i nawracania. Idiomatyczny Prolog używa głów klauzul do
analizy przypadków, akumulatorów do iteracji oraz porażki (a nie wyjątków) do wyrażenia „braku
odpowiedzi”. Odcięć i stanu globalnego używa się oszczędnie.

\section{Benchmarki generowania kodu}
\label{sec:related-benchmarks}

\paragraph{Benchmarki na poziomie funkcji} HumanEval~\cite{chen2021codex} zawiera 164 ręcznie
napisane problemy w Pythonie, z których każdy składa się z sygnatury funkcji, docstringu i średnio
7,7 testu jednostkowego. Wprowadził metrykę pass@$k$. MBPP~\cite{austin2021mbpp} zawiera 974 krótkie
problemy w Pythonie zebrane metodą crowdsourcingu. Oba stały się domyślnymi miarami dla modeli
kodu. \citet{liu2023evalplus} pokazali, że ich testy są często zbyt słabe: rozszerzenie HumanEval
o wiele dodatkowych testów (HumanEval+) wyraźnie obniża raportowane wyniki, ponieważ rozwiązania
wyglądające wiarygodnie, lecz błędne, przestają być akceptowane.

\paragraph{Benchmarki wielojęzyczne} MultiPL-E~\cite{cassano2023multiple} tłumaczy problemy i testy
z HumanEval i MBPP na 18 kolejnych języków za pomocą kompilatorów z Pythona do każdego z języków
docelowych. HumanEval-X~\cite{zheng2023codegeex} i MBXP~\cite{athiwaratkun2023mbxp} realizują tę
samą ideę. Benchmarki te pokazały, że wyniki spadają wraz z ilością danych treningowych dostępnych
dla danego języka~\cite{cassano2023multiple,cassano2024knowledge}. Problemy pozostają jednak
przetłumaczonymi problemami z Pythona. Są to krótkie, bezstanowe funkcje na listach, napisach
i liczbach, które rzadko wymagają projektowania specyficznego dla paradygmatu. Programowanie
logiczne jest w nich niemal nieobecne.

\paragraph{Większe jednostki i kontaminacja} ClassEval~\cite{du2023classeval} ocenia generowanie
całych klas, a SWE-bench~\cite{jimenez2024swebench} rozwiązywanie rzeczywistych zgłoszeń z GitHuba.
LiveCodeBench~\cite{jain2024livecodebench} z czasem dobiera nowe problemy z konkursów, aby uniknąć
kontaminacji, czyli obecności problemów z benchmarku i ich rozwiązań w danych treningowych, co
zawyża wyniki popularnych benchmarków~\cite{riddell2024contamination}.

Nasz benchmark sytuuje się między tymi nurtami. Jego zadania są większe niż problemy HumanEval
(kontrakty z wieloma funkcjami, struktury danych, własność pamięci), ale wystarczająco małe, by
zaimplementować je od zera w siedmiu językach z \emph{identycznymi} testami. Według naszej wiedzy
żaden istniejący benchmark nie utrzymuje jednocześnie stałego zadania i stałych testów we wszystkich
czterech klasycznych paradygmatach ani nie ocenia przy tym idiomatyczności wyniku.

\section{Samonaprawa z informacją zwrotną z wykonania}
\label{sec:related-repair}

Gdy wygenerowany kod zawodzi, można pokazać modelowi błąd i poprosić o poprawkę.
Self-Debugging~\cite{chen2024selfdebug} uczy modele wyjaśniać i naprawiać swój kod na podstawie
wyników wykonania, a Reflexion~\cite{shinn2023reflexion} pozwala agentowi przechowywać werbalną
informację zwrotną między próbami. \citet{olausson2024selfrepair} stwierdzili, że korzyść
z samonaprawy jest rzeczywista, lecz mniejsza, niż się często zakłada. Ogranicza ją przede wszystkim
jakość informacji zwrotnej, a tę kontroluje właśnie program uruchamiający testy. Nasz system
realizuje samonaprawę z informacją zwrotną z wykonania w ramach jednej konwersacji. Model widzi
komunikaty kompilatora, logi testów i nazwy niezaliczonych testów, ale nigdy kodu testów
(\cref{sec:repair-loop}).

\section{LLM jako sędziowie}
\label{sec:related-judge}

Użycie silnego modelu LLM do oceniania wyników (ang.~\emph{LLM-as-a-judge}) skaluje ewaluację tam,
gdzie nie istnieje odpowiedź wzorcowa. \citet{zheng2023judging} raportują zgodność między silnymi
sędziami LLM a ludźmi porównywalną ze zgodnością między ludźmi. Dokumentują też obciążenia związane
z pozycją, rozwlekłością i preferowaniem własnych odpowiedzi. \citet{panickssery2024selfpreference}
pokazują, że ewaluatory LLM potrafią rozpoznawać i faworyzować własne generacje. W przypadku kodu
ICE-Score~\cite{zhuo2024icescore} używa LLM do oceny użyteczności i poprawności funkcjonalnej bez
wykonywania kodu. My używamy sędziego wyłącznie do tego, czego testy nie mogą zmierzyć, czyli
idiomatyczności, czytelności i projektu, i ograniczamy go rubryką właściwą dla języka,
ustrukturyzowaną odpowiedzią w formacie JSON i zaślepieniem (\cref{sec:judge}).

\section{Języki programowania a jakość kodu}
\label{sec:related-quality}

Spór o to, czy wybór języka wpływa na jakość oprogramowania, toczył się na długo przed pojawieniem
się LLM. \citet{ray2014largescale} raportowali niewielkie, lecz istotne związki między językami
a commitami naprawiającymi defekty na GitHubie. Późniejsza replikacja \citet{berger2019impact}
wykazała, że większość tych efektów znika po skorygowaniu danych i analizy statystycznej. Ta historia
skłania do ostrożności: efekty języka łatwo pomylić z efektami projektu, danych i pomiaru, a ich
badanie wymaga starannej statystyki. Wracamy do tej kwestii w \cref{sec:threats}.

% =============================================================================
\chapter{Eksperymenty wstępne i wnioski}
\label{ch:preliminary}
% =============================================================================

Benchmark i system ewaluacyjny opisane w kolejnych rozdziałach nie były naszym pierwszym projektem.
Dwa kierunki pilotażowe zakończyły się niepowodzeniem, a przyczyny tych niepowodzeń ukształtowały
ostateczną metodykę. Opisujemy je, ponieważ są to przydatne wyniki negatywne dla każdego, kto planuje
podobne badanie.

\section{Zadania w stylu HumanEval są zbyt łatwe}
\label{sec:humaneval}

\paragraph{Motywacja} Naturalnym punktem wyjścia był HumanEval~\cite{chen2021codex}. Jest to
najczęściej raportowany benchmark generowania kodu, jego problemy są z założenia niezależne od
języka, a MultiPL-E~\cite{cassano2023multiple} pokazał, jak przetłumaczyć je wraz z testami na inne
języki. Planowaliśmy przenieść podzbiór jego problemów do siedmiu języków badania i porównać odsetki
rozwiązań.

\paragraph{Obserwacja} Problemy okazały się zdecydowanie zbyt łatwe. W naszym pilotażu modele
rozwiązały praktycznie każdy problem za pierwszym podejściem w każdym języku, w tym w językach
o małych zasobach. Wyniki wszystkich języków znalazły się na pułapie lub tuż pod nim, więc nie mogły
wykazać żadnych różnic między paradygmatami. Nie dało się też badać samonaprawy ani
idiomatyczności, ponieważ nie było czego naprawiać, a rozwiązania były zbyt krótkie, by ujawnić
jakikolwiek projekt.
\todo{dodać liczby z pilotażu: modele, liczba problemów i języków, odsetki rozwiązań}

\paragraph{Przyczyny} Z perspektywy czasu efekt sufitowy był do przewidzenia, z czterech powodów.
\begin{enumerate}
  \item \emph{Rozmiar problemu.} Typowy problem HumanEval wymaga jednej krótkiej funkcji na
    listach, napisach lub liczbach całkowitych. Taka funkcja ma kilka linii w dowolnym języku i nie
    wymaga ani struktur danych, ani zarządzania własnością pamięci, więc nie zostawia miejsca na
    dekompozycję specyficzną dla paradygmatu.
  \item \emph{Nasycenie.} Obecne czołowe modele rozwiązują zdecydowaną większość problemów
    HumanEval, dlatego społeczność przeszła do trudniejszych lub nowszych
    benchmarków~\cite{liu2023evalplus,jain2024livecodebench,jimenez2024swebench}.
  \item \emph{Kontaminacja.} Problemy i wiele rozwiązań są publiczne od 2021 roku i prawdopodobnie
    znajdują się w danych treningowych obecnych modeli~\cite{riddell2024contamination}. Wysoki wynik
    może wtedy świadczyć o zapamiętaniu, a nie o umiejętności programowania.
  \item \emph{Słabe testy.} Przy niewielkiej liczbie testów na problem przechodzą rozwiązania
    częściowo błędne~\cite{liu2023evalplus}. Przez to pułap wydaje się jeszcze wyższy, niż jest
    w rzeczywistości.
\end{enumerate}

\paragraph{Konsekwencje dla projektu} Postanowiliśmy napisać własne zadania, spełniające poniższe
wymagania, z których każde odpowiada na jeden z powyższych problemów:
\begin{description}[leftmargin=2.4em,labelwidth=2em,itemsep=2pt]
  \item[D1] \emph{Kontrakty z wieloma funkcjami.} Każde zadanie wymaga dwóch powiązanych algorytmów
    za niewielkim API z własnymi typami danych, np.\ grafem, drzewem AST lub B-drzewem, tak aby
    rozwiązanie trzeba było zaprojektować, a nie tylko napisać.
  \item[D2] \emph{Precyzyjne specyfikacje.} Reguły rozstrzygania remisów, wartości błędów
    i przypadki brzegowe są wyspecyfikowane tak, aby oczekiwany wynik był jednoznaczny i dało się
    go porównać dokładnie (na przykład dokładny kształt B-drzewa po ciągu wstawień i usunięć).
  \item[D3] \emph{Mocne testy.} Każde zadanie ma 20 testów (10 na algorytm), w tym dane puste
    i zdegenerowane, niepoprawne argumenty oraz testy obciążeniowe na dużych danych.
  \item[D4] \emph{Ukryte testy.} Model nigdy nie widzi kodu testów, więc nie może dopasować
    rozwiązania do testów.
  \item[D5] \emph{Bezpieczeństwo pamięci.} Tam, gdzie język na to pozwala, testy są uruchamiane pod
    Valgrindem, więc wycieki i niepoprawne dostępy do pamięci liczą się jako błędy.
  \item[D6] \emph{Kompozycje trudniejsze do zapamiętania.} Klasyczne algorytmy są łączone
    z nietypowymi wymaganiami, takimi jak ciągłe wykrywanie kolizji oparte na sumie Minkowskiego,
    selekcja BFPRT raportująca ślad wyboru pivotów czy mediana dwóch posortowanych ciągów przy
    ścisłym limicie odczytów, aby zmniejszyć wartość zapamiętanych rozwiązań.
\end{description}
\Cref{ch:benchmark} opisuje powstały benchmark. \Cref{ch:results} pokazuje, że unika on efektu
sufitowego: za pierwszym podejściem rozwiązano tylko 80\,\% zadań, a błędy różnią się między
językami.

\section{Lokalne modele o otwartych wagach są niepraktyczne}
\label{sec:local-llms}

\paragraph{Motywacja} Drugim planem było lokalne uruchamianie modeli o otwartych wagach. Obiecywało
to zerowy koszt krańcowy przebiegu, pełną kontrolę nad dekodowaniem (temperatura, ziarna losowe),
stałe wagi modelu zapewniające powtarzalność oraz niezależność od komercyjnego API.
\todo{podać nazwy modeli lokalnych (rodzina, rozmiar, kwantyzacja), środowisko inferencji i sprzęt}

\paragraph{Przeszkoda 1: wydajność} Generowanie na dostępnym nam sprzęcie było zbyt wolne jak na
badanie tej skali. Pojedyncza próba składa się z długiego promptu (specyfikacja wraz z kontraktem,
a w rundach naprawy także poprzedni kod i log testów) i długiej odpowiedzi (kompletny moduł
o długości 100--280 linii, zob.\ \cref{tab:efficiency}). Konwersacja naprawcza rośnie z każdą
rundą. Dla porównania: w naszym ostatecznym przebiegu \haiku{} generował średnio około 6\,500
tokenów wyjściowych na wywołanie, łącznie z rozumowaniem, i 635\,269 tokenów wyjściowych w całym
przebiegu. Przy przykładowej lokalnej szybkości dekodowania 10~tokenów/s zajęłoby to około
11~minut na wywołanie i prawie 18~godzin samego dekodowania na jedno przejście przez benchmark, bez
przetwarzania promptów i bez powtórzonych próbek, których wymaga badanie statystyczne. Mniejsze lub
silniej skwantyzowane modele, które działały szybciej, pisały wyraźnie gorszy kod, więc mierzyłyby
ograniczenia sprzętu, a nie wpływ paradygmatu.
\todo{dodać zmierzoną przepustowość (tokeny/s) i czas rzeczywisty na zadanie z pilotażu}

\paragraph{Przeszkoda 2: modele nie zwracały samego kodu} System zapisuje odpowiedź modelu
bezpośrednio do pliku źródłowego, kompiluje go i uruchamia testy. Wymaga to, aby odpowiedź była
kodem i niczym więcej. Modele lokalne często się do tego nie stosowały, nawet gdy prompt mówił o tym
wprost. Dodawały wyjaśnienia przed kodem i po nim, otaczały go znacznikami bloków Markdown, podawały
kilka alternatywnych wersji, zwracały fragmenty zamiast kompletnych modułów albo powtarzały części
promptu. Każda taka odpowiedź kończyła się błędem kompilacji, który mierzył wykonywanie instrukcji,
a nie umiejętność programowania. Co gorsza, efekt ten nie był jednakowy dla wszystkich języków,
więc mieszał się ze zmienną, którą chcieliśmy badać. Heurystyki wycinające kod z prozy (wybór
największego bloku kodu, usunięcie tekstu wstępnego) rozwiązywały proste przypadki. Zawodziły,
gdy bloków kodu było kilka lub brakowało znaczników, a każda ręczna decyzja, który blok jest
„odpowiedzią”, czyniłaby wyniki mniej obiektywnymi. Ostatecznie wyniki tych przebiegów były bardzo
trudne do interpretacji i zrezygnowaliśmy w tym badaniu z modeli lokalnych.

\paragraph{Decyzja} Przeszliśmy na hostowane modele Claude dostępne przez API Anthropic. Ten wybór
usuwa obie przeszkody. Przepustowość jest wystarczająca, by przeprowadzić cały benchmark w około
trzy kwadranse czasu pracy modelu (\cref{sec:cost}). Modele oferują też \emph{adaptacyjne
myślenie} (ang.~\emph{adaptive thinking})~\cite{anthropic2026docs}, które zwraca rozumowanie
w osobnych blokach \emph{thinking}, dzięki czemu odpowiedź tekstowa może być czystym kodem. Sędzia
korzysta z \emph{ustrukturyzowanych odpowiedzi} (ang.~\emph{structured outputs}), tj.\ schematu
JSON, który odpowiedź musi spełniać. Zachowaliśmy awaryjny mechanizm wycinania kodu i każdą
odpowiedź, która go wymagała, rejestrujemy jako \emph{naruszenie formatu}. Dzięki temu
wykonywanie instrukcji staje się mierzoną wielkością, a nie ukrytym źródłem szumu
(\cref{sec:extraction}).

\begin{findingbox}[Wniosek.]
Zgodność z wymaganym formatem odpowiedzi jest warunkiem wstępnym zautomatycznej ewaluacji kodu, a nie
szczegółem. Nawet model hostowany otoczył kod znacznikami bloków Markdown w 37,8\,\% odpowiedzi,
a w dwóch odpowiedziach wplótł do pliku źródłowego tekst niebędący kodem
(\cref{sec:cost,sec:failure-analysis}). Benchmark powinien zatem mierzyć i raportować zgodność
formatu oddzielnie od poprawności.
\end{findingbox}

% =============================================================================
\chapter{Projekt benchmarku}
\label{ch:benchmark}
% =============================================================================

Ten rozdział opisuje benchmark: języki i ich środowiska wykonawcze, strukturę zadania i samych
dziesięć zadań oraz sposób walidacji zestawów testów.

\section{Języki i środowiska wykonawcze}
\label{sec:languages}

\Cref{tab:languages} zestawia siedem języków. Katalog każdego języka zawiera \file{Dockerfile}
z narzędziami, \file{Makefile} i skrypt \file{run-tests.sh}, który buduje i uruchamia katalog jednego
zadania, oraz plik \file{README.md} z dwoma szablonami promptów (\cref{sec:prompts}). Testy są
uruchamiane w jednorazowych kontenerach, więc wadliwe rozwiązanie nie może wpłynąć na system
gospodarza ani na inne przebiegi. W C, Pascalu, Haskellu i Prologu program testowy działa pod
Valgrindem~\cite{nethercote2007valgrind} z opcją \texttt{-{}-error-exitcode=1}. Rozwiązanie, które
przechodzi wszystkie testy, ale ma wycieki pamięci lub odczytuje niepoprawną pamięć, jest więc mimo
to uznawane za błędne. W Prologu sprawdzanie wycieków jest wyłączone, ponieważ alokacje samego
interpretera powodowałyby niepowodzenie każdego przebiegu.

\begin{table}[htbp]
  \centering
  \caption{Języki, paradygmaty i środowiska wykonawcze. \emph{Kontrakt} to plik pokazywany
    modelowi, a \emph{cel} to plik zastępowany jego odpowiedzią; „moduł” oznacza, że model
    otrzymuje moduł-szkielet i zwraca go w kompletnej postaci.}
  \label{tab:languages}
  \small
  \begin{tabularx}{\textwidth}{@{}llYllc@{}}
    \toprule
    Język & Paradygmat & Narzędzia (bazowy obraz Docker) & Kontrakt & Cel & Valgrind \\
    \midrule
    C         & imperatywny & GCC 12, Debian~12 (\code{debian:bookworm-slim})    & \file{name.h}    & \file{name.c}        & wycieki + pamięć \\
    Pascal    & imperatywny & Free Pascal 3.2.2, Debian~12                        & \multicolumn{2}{l}{moduł (\file{.pas})} & wycieki + pamięć \\
    Java      & obiektowy   & JDK 17, Eclipse Temurin (\code{eclipse-temurin:17-jdk})  & \file{Name.java} & \file{NameImpl.java} & -- \\
    Smalltalk & obiektowy   & GNU Smalltalk 3.2.5, Debian~11                      & \file{Name.st}   & \file{NameImpl.st}   & -- \\
    Haskell   & funkcyjny   & GHC 9.0.2, Debian~12                                & \multicolumn{2}{l}{moduł (\file{.hs})}  & wycieki + pamięć \\
    Scala     & funkcyjny   & Scala 2.11.12 na domyślnym JDK, Debian~12          & \multicolumn{2}{l}{moduł (\file{.scala})} & -- \\
    Prolog    & logiczny    & SWI-Prolog 9.2.9 (\code{swipl:9.2.9})               & \multicolumn{2}{l}{moduł (\file{.pl})}  & tylko pamięć \\
    \bottomrule
  \end{tabularx}
\end{table}

\section{Struktura zadania}
\label{sec:task-structure}

Zadanie to katalog \file{<Language>/test-N/}. To samo zadanie istnieje w każdym języku, z tymi
samymi testami, danymi wejściowymi i oczekiwanymi wynikami. Różni się tylko reprezentacja, np.\
\code{NULL} w C, \code{Nothing} w Haskellu i atom \code{none} w Prologu. \Cref{tab:task-files}
wymienia pliki katalogu zadania.

\begin{table}[htbp]
  \centering
  \caption{Pliki katalogu zadania i ich widoczność dla modelu.}
  \label{tab:task-files}
  \small
  \begin{tabularx}{\textwidth}{@{}lYl@{}}
    \toprule
    Plik & Rola & Widoczny dla modelu \\
    \midrule
    \file{test-N.md}            & Specyfikacja zadania; blok pod nagłówkiem \texttt{\#\# Prompt} jest wstawiany do promptu & tak \\
    kontrakt (\file{knapsack.h}, \file{Knapsack.java}, \ldots) & Typy i sygnatury, które rozwiązanie musi zaimplementować & tak \\
    szkielet (\file{knapsack.c}, \file{KnapsackImpl.java}, \ldots) & Ciała typu „nie zaimplementowano”; zastępowane odpowiedzią modelu & jako kontrakt w językach jednoplikowych \\
    testy (\file{tests.c}, \file{Tests.hs}, \ldots) & Testy jednostkowe i program je uruchamiający & \textbf{nigdy} \\
    \bottomrule
  \end{tabularx}
\end{table}

\paragraph{Dwa algorytmy, dwadzieścia testów} Każde zadanie wymaga dwóch powiązanych algorytmów,
A i B, a każdy algorytm ma dziesięć testów nazwanych od \code{Test A01 (Title)} do
\code{Test B10 (Title)}. System odczytuje identyfikatory testów z pliku testowego i nie zakłada
stałej liczby testów. Testy obejmują przypadki typowe, przypadki brzegowe (puste dane, pojedynczy
element, argumenty \code{NULL}, niepoprawne indeksy), podchwytliwe przypadki specyficzne dla
algorytmu oraz co najmniej jeden test obciążeniowy na dużych danych.

\paragraph{Protokół wyjścia} Wszystkie programy testowe wypisują ten sam protokół
(\cref{lst:protocol}) i kończą się kodem~0 tylko wtedy, gdy wszystkie testy przeszły, a Valgrind,
jeśli jest używany, niczego nie wykrył. Znaczniki \code{=== START ===} i
\code{=== BENCHMARK RESULTS ===} pozwalają systemowi odróżnić kod, który się nie skompilował, od
kodu, który uległ awarii w trakcie testów (\cref{sec:outcomes}).

\begin{lstlisting}[style=prompt,caption={Protokół wyjścia testów wspólny dla wszystkich języków (fragment przebiegu zadania Dijkstry w Smalltalku).},label={lst:protocol}]
=== START ===
--- Algorithm A: Shortest Distances ---
[FAIL] Test A01 (Linear Chain)
...
[PASS] Test A07 (Null Graph)
--- Algorithm B: Path Reconstruction ---
[PASS] Test B01 (Simple Path)
...
=== BENCHMARK RESULTS ===
\end{lstlisting}

\paragraph{Kontrakty w różnych paradygmatach} Kontrakt ustala API, ale pozostawia implementację
otwartą. \Cref{lst:contract-c,lst:contract-hs,lst:contract-pl} pokazują kontrakt zadania plecakowego
w trzech paradygmatach. W C model implementuje plik nagłówkowy i musi zwracać wyniki zaalokowane na
stercie, które zwalnia wywołujący. W Haskellu uzupełnia moduł z algebraicznym typem danych
i czystymi funkcjami. W Prologu implementuje moduł z relacjami, w którym wynikami są termy takie jak
\code{result(MaxValue, Selected, TotalWeight)}. W językach jednoplikowych (Pascal, Haskell, Scala,
Prolog) model zwraca cały moduł wraz z deklaracjami typów, więc może też naruszyć kontrakt; objawia
się to błędem kompilacji.

\begin{lstlisting}[language=C,caption={Kontrakt zadania 4 (plecak 0/1) w C: \file{knapsack.h}.},label={lst:contract-c}]
typedef struct {
    int max_value;
    int* selected;       /* length n; entries 0 or 1 */
    int n;
    int total_weight;
} KnapsackResult;

int knapsack_max_value(const int* weights, const int* values, int n, int capacity);
KnapsackResult* knapsack_solve(const int* weights, const int* values, int n, int capacity);
void free_knapsack_result(KnapsackResult* result);
\end{lstlisting}

\begin{lstlisting}[language=Haskell,caption={Ten sam kontrakt w Haskellu: moduł-szkielet \file{Knapsack.hs}.},label={lst:contract-hs}]
module Knapsack ( Result (..), maxValue, solve ) where

data Result = Result
  { resultMaxValue    :: Int
  , resultSelected    :: [Int]
  , resultTotalWeight :: Int
  } deriving (Eq, Show)

maxValue :: [Int] -> [Int] -> Int -> Int
maxValue _ _ _ = error "Knapsack.maxValue not implemented"

solve :: [Int] -> [Int] -> Int -> Result
solve _ _ _ = error "Knapsack.solve not implemented"
\end{lstlisting}

\begin{lstlisting}[language=Prolog,caption={Ten sam kontrakt w Prologu: moduł-szkielet \file{knapsack.pl} z trybami argumentów udokumentowanymi w komentarzach.},label={lst:contract-pl}]
:- module(knapsack, [knapsack_max_value/4, knapsack_solve/4]).

% knapsack_max_value(+Weights, +Values, +Capacity, -MaxValue)
% knapsack_solve(+Weights, +Values, +Capacity, -Result)
%   Result = result(MaxValue, Selected, TotalWeight)

knapsack_max_value(_, _, _, _) :-
    throw(error(existence_error(procedure, knapsack_max_value/4), _)).
knapsack_solve(_, _, _, _) :-
    throw(error(existence_error(procedure, knapsack_solve/4), _)).
\end{lstlisting}

\section{Dziesięć zadań}
\label{sec:tasks}

\Cref{tab:tasks} zestawia zadania. Obejmują one geometrię obliczeniową, grafy, struktury danych na
napisach, programowanie dynamiczne, parsowanie, przeszukiwanie z ograniczeniami, bioinformatykę,
algorytmy aproksymacyjne, struktury danych dla pamięci zewnętrznej i selekcję. Nowsze zadania (7--10)
są wyspecyfikowane ściślej niż starsze: definiują każdą regułę rozstrzygania remisów, tak aby
oczekiwany wynik był jednoznaczny i porównywany dokładnie.

\begin{table}[htbp]
  \centering
  \caption{Dziesięć zadań benchmarku. Każdy algorytm ma dziesięć testów (20 na zadanie, 200 na
    język, łącznie 1\,400).}
  \label{tab:tasks}
  \small
  \begin{tabularx}{\textwidth}{@{}rL{2.7cm}L{3.0cm}L{3.0cm}Y@{}}
    \toprule
    \# & Zadanie & Algorytm A & Algorytm B & Co stanowi trudność \\
    \midrule
    1  & Logika przestrzenna        & Oczyszczona suma Minkowskiego (otoczka wypukła) & Dynamiczne wykrywanie kolizji & Wielokąty zdegenerowane (punkty, odcinki, puste), usuwanie wierzchołków współliniowych, kolejność przeciwna do ruchu wskazówek zegara, ciągła kolizja w przestrzeni konfiguracyjnej~\cite{deberg2008cg,andrew1979hull} \\
    2  & Dijkstra                   & Najkrótsze odległości & Odtwarzanie ścieżki & Kolejka priorytetowa, krawędzie równoległe i o wadze zero, wierzchołki nieosiągalne ($\infty$), łańcuchy o 1\,000 wierzchołków~\cite{dijkstra1959note} \\
    3  & Eertree                    & Budowa i zapytania strukturalne & Zliczanie wystąpień & Drzewo palindromowe w czasie liniowym z dwoma korzeniami, dowiązaniami sufiksowymi i propagacją liczników; test obciążeniowy na 50\,000 znaków~\cite{rubinchik2018eertree} \\
    4  & Plecak 0/1                 & Maksymalna wartość & Wybór przedmiotów & Programowanie dynamiczne z odtworzeniem optymalnego wyboru; pułapka na algorytm zachłanny~\cite{kellerer2004knapsack} \\
    5  & Parser wyrażeń             & Parsowanie i budowa AST & Wartościowanie i semantyka błędów & Zejście rekurencyjne z priorytetami i łącznością ($-2^2$, $2^{3^2}$), arność funkcji, propagacja pierwszego statusu błędu~\cite{aho2006compilers} \\
    6  & Sudoku                     & Walidacja i kandydaci & Rozwiązywanie i zliczanie rozwiązań & Nawracanie z heurystyką MRV, zliczanie z wczesnym zatrzymaniem; naiwne przeszukiwanie jest zbyt wolne~\cite{russell2020aima} \\
    7  & Dopasowanie sekwencji      & Wynik dopasowania & Odtworzenie dopasowania & Programowanie dynamiczne Gotoha z afinicznymi karami za przerwy, trzema macierzami i przejściami przerwa--przerwa; ścieżka powrotna między macierzami~\cite{gotoh1982improved} \\
    8  & Aproksymacja metrycznego TSP & Podwojone drzewo MST & Heurystyka Christofidesa & MST algorytmem Prima z dokładnymi regułami remisów, dokładne minimalne skojarzenie doskonałe (DP na maskach bitowych), cykl Eulera algorytmem Hierholzera, skróty~\cite{rosenkrantz1977tsp,christofides1976worst,prim1957shortest} \\
    9  & B-drzewo                   & Wstawianie & Usuwanie & Podział i scalanie węzłów z góry na dół według CLRS; porównywany jest dokładny kształt drzewa~\cite{bayer1972btree,cormen2009clrs} \\
    10 & Statystyki pozycyjne       & Mediana median ze śladem pivotów & Mediana dwóch posortowanych źródeł & Selekcja BFPRT z dokładnym śladem pivotów; wyszukiwanie binarne przy limicie $16(\lfloor\log_2(m{+}n)\rfloor+2)$ odczytów~\cite{blum1973time} \\
    \bottomrule
  \end{tabularx}
\end{table}

\noindent Zadania bardziej szczegółowo:

\begin{description}[leftmargin=0pt,itemsep=3pt,font=\normalfont\bfseries]
  \item[Zadanie 1: logika przestrzenna.] Oblicz sumę Minkowskiego dwóch wielokątów wypukłych i zwróć
    tylko ściśle skrajne punkty jej otoczki wypukłej w kolejności przeciwnej do ruchu wskazówek
    zegara. Następnie rozstrzygnij, czy wielokąt robota poruszający się ze stałą prędkością zderzy
    się w dowolnym miejscu swojej trajektorii z nieruchomą przeszkodą, korzystając z przeszkody
    w przestrzeni konfiguracyjnej, aby uniknąć „tunelowania”. Wersja w C zarządza też strukturami
    wielokątów i warstw zawierającymi posiadane napisy.
  \item[Zadanie 2: Dijkstra.] Zbuduj skierowany graf ważony z listami sąsiedztwa, oblicz najkrótsze
    odległości z użyciem kopca minimalnego w czasie $O((V+E)\log V)$ i odtwórz ścieżki z tablicy
    poprzedników. Niepoprawne wierzchołki źródłowe i docelowe mają zdefiniowane wyniki
    (\code{NULL}, \code{INFINITY}, \code{none}).
  \item[Zadanie 3: eertree.] Zbuduj drzewo palindromowe napisu w czasie liniowym, policz różne
    podsłowa będące palindromami, znajdź najdłuższe z nich i policz wystąpienia danego palindromu
    po propagacji liczników wzdłuż dowiązań sufiksowych.
  \item[Zadanie 4: plecak 0/1.] Zwróć optymalną wartość i realizujący ją wybór przedmiotów,
    z zachowaniem niezmienników dotyczących wagi i wartości wyboru. Specyfikacja wprost ostrzega,
    że heurystyki zachłanne zawodzą.
  \item[Zadanie 5: parser wyrażeń.] Sparsuj wyrażenia arytmetyczne ze zmiennymi i siedmioma
    funkcjami wbudowanymi do drzewa AST zgodnie z podaną gramatyką. Oblicz jego wartość
    z typowanymi statusami błędów (niezdefiniowana zmienna, dzielenie przez zero, błąd dziedziny),
    które przerywają obliczenie.
  \item[Zadanie 6: Sudoku.] Waliduj plansze, wypisz kandydatów dla komórki, rozwiąż łamigłówkę
    i policz rozwiązania do zadanego limitu. Nawracanie musi rozgałęziać się na komórce z najmniejszą
    liczbą kandydatów (MRV), ponieważ naiwne przeszukiwanie wierszami jest zbyt wolne dla części
    testowych łamigłówek.
  \item[Zadanie 7: dopasowanie sekwencji.] Oblicz wynik optymalnego dopasowania globalnego
    z afinicznymi karami za przerwy (przerwa długości $L$ kosztuje $o + L\cdot e$) i zwróć jedno
    optymalne dopasowanie. Dane o długości ponad 100 znaków wykluczają pełne wyliczanie.
  \item[Zadanie 8: metryczny TSP.] Dla odległości Manhattan oblicz wagę MST i trasę podwojonego
    drzewa (jednoznaczną dzięki podanym regułom remisów), dokładne skojarzenie doskonałe o minimalnej
    wadze na wierzchołkach stopnia nieparzystego oraz trasę Christofidesa poprzez cykl Eulera
    w multigrafie, z zachowaniem krawędzi podwójnych.
  \item[Zadanie 9: B-drzewo.] Wstawiaj i usuwaj klucze całkowite zgodnie z CLRS~\cite{cormen2009clrs},
    z wyprzedzającym dzieleniem i scalaniem węzłów, łącznie ze sprawdzeniem przynależności przed
    modyfikacją. Testy porównują dokładny kształt otrzymanego drzewa.
  \item[Zadanie 10: statystyki pozycyjne.] Zaimplementuj selekcję BFPRT~\cite{blum1973time} na
    grupach po pięć elementów ze stabilnym podziałem trójdzielnym i raportuj dokładny ciąg pivotów.
    Oblicz medianę dwóch posortowanych ciągów, które można odczytywać wyłącznie przez zliczającą
    sondę; przekroczenie logarytmicznego limitu odczytów przerywa wywołanie.
\end{description}

\section{Walidacja benchmarku}
\label{sec:benchmark-validation}

Błędny test karałby każdy model i każdy język, który rozwiązuje zadanie poprawnie, dlatego zestawy
testów są walidowane przed każdym płatnym przebiegiem.

\paragraph{Tryb bazowy} Opcja \texttt{-{}-baseline} systemu przepuszcza przez ewaluator każdy szkielet
i każde dostępne rozwiązanie wzorcowe, bez wywoływania modelu. Każdy szkielet musi się
skompilować, wypisać linię wyniku dla każdego testu i nie przejść testów. Rozwiązania wzorcowe,
przechowywane w historii gita, muszą przechodzić tam, gdzie się tego oczekuje. Istnieje pięć
rozwiązań wzorcowych (C: zadania 1 i 3, Pascal: zadanie 1, Java: zadanie 1, Smalltalk: zadanie 1).
Rozwiązania w Javie i Smalltalku nie przechodzą kilku testów i nie są używane jako dowód
poprawności testów.

\paragraph{Testy przechodzące na szkielecie} Niektóre testy oczekują dokładnie tego, co zwraca
szkielet, np.\ \code{NULL} lub \code{false} dla niepoprawnych danych. Tryb bazowy wykrył 17 takich
testów, wszystkie w C (sześć w zadaniu~1, pięć w zadaniu~5, po dwa w zadaniach~2 i~10, po jednym
w zadaniach~3 i~9). Nieznacznie zawyżają one liczbę zaliczonych testów, dlatego raport podaje wyniki
dla C z nimi i bez nich (\cref{sec:results-correctness}).

\paragraph{Generowane zestawy testów} Zadania 1--5 zostały napisane ręcznie, język po języku.
Zadania 6--10 wytworzono za pomocą umiejętności (ang.~\emph{skill}) Claude Code
(\file{/new-benchmark}) napisanej na potrzeby tego projektu. Umiejętność najpierw projektuje
niezależną od języka specyfikację z dwoma algorytmami i dziesięcioma testami dla każdego z nich.
Po zatwierdzeniu przez autorów generuje kontrakt, szkielet, testy i specyfikację dla każdego języka
i weryfikuje je względem tymczasowego rozwiązania wzorcowego, które następnie jest usuwane. Nowsze
zadania są więc bardziej jednolite między językami. Jak pokazuje \cref{sec:benchmark-defects}, oba
defekty ujawnione przez przebieg znajdują się w ręcznie napisanych zestawach dla Prologu.

% =============================================================================
\chapter{System ewaluacyjny}
\label{ch:harness}
% =============================================================================

System ewaluacyjny (ang.~\emph{harness}) to program w Pythonie (około 1\,900 linii, \file{harness/}),
który przeprowadza całą ewaluację: generuje prompty, wywołuje model, ocenia odpowiedzi w Dockerze,
steruje pętlą naprawy, wywołuje sędziego i zapisuje raport. Ten rozdział opisuje jego architekturę,
a następnie każdy etap, ze szczególnym uwzględnieniem konstrukcji promptów.

\section{Architektura}
\label{sec:architecture}

\Cref{fig:pipeline} przedstawia potok przetwarzania, a \cref{tab:modules} moduły. Wszystkie
ustawienia znajdują się w jednym pliku, \file{config.py}, a flagi wiersza poleceń nadpisują je dla
pojedynczego przebiegu. Program ma trzy tryby: \texttt{-{}-dry-run} generuje wszystkie prompty bez
wywoływania API i Dockera, \texttt{-{}-baseline} waliduje zestawy testów
(\cref{sec:benchmark-validation}), a tryb domyślny wykonuje płatny przebieg po wyświetleniu
szacunku kosztów i prośbie o potwierdzenie.

\begin{figure}[htbp]
  \centering
  % Pipeline of the evaluation harness (used in chapters/05-harness.tex)
\begin{tikzpicture}[
  font=\footnotesize,
  node distance=12mm and 7mm,
  box/.style={draw=inksec, rounded corners=2pt, fill=white, align=center,
              minimum height=11mm, text width=25mm, inner sep=3pt},
  llm/.style={box, draw=seq600, fill=seq050, line width=0.7pt},
  store/.style={box, draw=inkmuted, fill=neutralbg},
  dec/.style={draw=inksec, diamond, aspect=2, fill=white, align=center, inner sep=0.5pt,
              font=\scriptsize, minimum width=22mm},
  arr/.style={-{Stealth[length=2.2mm]}, draw=inksec, line width=0.6pt},
  lbl/.style={font=\scriptsize, text=inksec, fill=white, inner sep=1pt},
]
  % row 1: task -> prompt -> generator -> extraction
  \node[store] (task) {Katalog zadania\\[-1pt] \scriptsize specyfikacja + kontrakt};
  \node[box, right=of task] (prompt) {Generator promptów\\[-1pt] \scriptsize szablony, symbole zastępcze};
  \node[llm, right=of prompt] (gen) {LLM-generator\\[-1pt] \scriptsize \haiku};
  \node[box, right=of gen] (extract) {Wycinanie kodu\\[-1pt] \scriptsize kontrola formatu};
  \node[store, right=of extract, text width=19mm] (tests) {Ukryte testy\\[-1pt] \scriptsize nigdy niepokazywane};

  % row 2 (flows right to left): evaluator -> classifier -> decision -> repair prompt
  \node[box, below=of extract] (docker) {Ewaluator Docker\\[-1pt] \scriptsize budowanie, testy, Valgrind};
  \node[box, below=of gen] (classify) {Klasyfikator wyniku\\[-1pt] \scriptsize 9 kategorii};
  \node[dec, below=of prompt] (passq) {rozwiązane?};
  \node[box, below=of task] (repair) {Prompt naprawczy\\[-1pt] \scriptsize fragment logu +\\[-2pt] niezaliczone testy};

  % row 3: judge and report
  \node[llm, below=of classify] (judge) {Sędzia LLM\\[-1pt] \scriptsize Sonnet 5.5, zaślepiony};
  \node[store] (final) at (passq |- judge) {Zapis próbki\\[-1pt] \scriptsize każda próba zapisana};
  \node[box, below=of docker] (report) {Raport\\[-1pt] \scriptsize tabele, CSV, koszty};

  % arrows
  \draw[arr] (task) -- (prompt);
  \draw[arr] (prompt) -- (gen);
  \draw[arr] (gen) -- (extract);
  \draw[arr] (extract) -- (docker);
  \draw[arr, dashed, draw=inkmuted] (tests.south) |- (docker.east);
  \draw[arr] (docker) -- (classify);
  \draw[arr] (classify) -- (passq);
  \draw[arr] (passq) -- node[lbl, above] {nie} (repair);
  \draw[arr] (repair.north) |- ($(prompt.south)+(-4mm,-5mm)$) -- ($(prompt.south)+(-4mm,0)$);
  \node[lbl, anchor=south west] at ($(repair.north)+(1mm,1.5mm)$) {$\le 3$ rundy, ta sama konwersacja};
  \draw[arr] (passq) -- node[lbl, right, align=left] {tak lub brak\\[-2pt]kolejnych rund} (final);
  \draw[arr] (final) -- (judge);
  \draw[arr] (judge) -- (report);
\end{tikzpicture}

  \caption{Potok systemu dla jednego zadania. Generator odpowiada w ramach jednej konwersacji; po
    nieudanej ewaluacji otrzymuje prompt naprawczy z fragmentem logu i nazwami niezaliczonych
    testów, przez co najwyżej trzy rundy. Kod testów nigdy nie jest pokazywany żadnemu z modeli.}
  \label{fig:pipeline}
\end{figure}

\begin{table}[htbp]
  \centering
  \caption{Moduły systemu ewaluacyjnego.}
  \label{tab:modules}
  \small
  \begin{tabularx}{\textwidth}{@{}lY@{}}
    \toprule
    Moduł & Odpowiedzialność \\
    \midrule
    \file{config.py}    & Wszystkie ustawienia: identyfikator przebiegu, modele, poziom wysiłku, rundy naprawy, liczba wątków, limity czasu, sędzia, ceny \\
    \file{run.py}       & Punkt wejścia i trzy tryby; pętla generowanie $\to$ ewaluacja $\to$ naprawa \\
    \file{tasks.py}     & Wyszukiwanie katalogów zadań, tabele języków i paradygmatów, pliki kontraktu i celu, identyfikatory testów \\
    \file{prompts.py}   & Generowanie promptów początkowych i naprawczych z szablonów \\
    \file{llm.py}       & Klient API, wywołanie generujące, wycinanie kodu, szacunek kosztów i potwierdzenie \\
    \file{evaluator.py} & Obrazy Docker, kopie robocze, uruchamianie testów, parsowanie wyjścia, klasyfikacja wyniku, fragmenty logów \\
    \file{judge.py}, \file{rubrics.md} & LLM jako sędzia idiomatyczności, rubryka, tryb kalibracji \\
    \file{tracing.py}   & Rejestrowanie każdego wywołania API (żądanie, odpowiedź, podsumowanie rozumowania, identyfikator żądania, czas, tokeny) \\
    \file{report.py}, \file{store.py} & Agregacja do \file{report.md}, \file{summary.csv} i plików JSONL; ścieżki wyników \\
    \bottomrule
  \end{tabularx}
\end{table}

\section{Konstrukcja promptów}
\label{sec:prompts}

\paragraph{Szablony} Katalog każdego języka zawiera plik \file{README.md} z dwoma szablonami
promptów: \emph{początkowym} i \emph{naprawczym} (ang.~\emph{follow-up}). System pobiera pierwszy blok
kodu pod nagłówkami \emph{Initial prompt} i \emph{Follow-up prompt} i wypełnia symbole zastępcze
z \cref{tab:placeholders}. Trzymanie szablonów obok narzędzi danego języka ułatwia ich przegląd.
Ich umiejscowienie wymusza też regułę: każda zmiana szablonu zmienia to, co widzi model, więc wyniki
sprzed zmiany i po niej muszą należeć do różnych przebiegów.

\begin{table}[htbp]
  \centering
  \caption{Symbole zastępcze szablonów promptów. Generowanie kończy się błędem, jeśli którykolwiek
    symbol pozostanie niewypełniony.}
  \label{tab:placeholders}
  \small
  \begin{tabularx}{\textwidth}{@{}lY@{}}
    \toprule
    Symbol zastępczy & Wypełniany przez \\
    \midrule
    \code{[TASK-SPECIFIC REQUIREMENTS]}      & blok pod nagłówkiem \texttt{\#\# Prompt} w \file{test-N.md} \\
    \code{[INSERT ... HERE]}                 & plik kontraktu w dosłownym brzmieniu (nagłówek, klasa abstrakcyjna, interfejs lub moduł-szkielet) \\
    \code{[ERROR LOGS ...]}                  & fragment logu nieudanej ewaluacji (tylko prompt naprawczy, \cref{sec:repair-loop}) \\
    \code{[FAILING TEST CASE ...]}           & nazwy niezaliczonych testów albo informacja, że żaden test nie dał wyniku (tylko prompt naprawczy) \\
    \code{[TARGET FILE NAME]}                & plik zapisywany przez model, np.\ \file{knapsack.c} (tylko prompt naprawczy) \\
    \bottomrule
  \end{tabularx}
\end{table}

\paragraph{Struktura promptu początkowego} Każdy prompt początkowy składa się z tych samych czterech
części (\cref{lst:prompt-initial}):
\begin{enumerate}
  \item zdania określającego \emph{rolę} („You are a C programming expert”) i cel (implementacja
    kontraktu);
  \item \emph{wymagań specyficznych dla zadania}, tj.\ specyfikacji obu algorytmów z przypadkami
    brzegowymi, wartościami błędów oraz, w razie potrzeby, wymaganiami dotyczącymi złożoności
    i algorytmu;
  \item \emph{wymagań ogólnych} danego języka: wskazówek dotyczących paradygmatu i stylu oraz
    formatu odpowiedzi (wyłącznie surowy kod, bez Markdownu, bez kodu testów);
  \item \emph{kontraktu}, dołączonego w dosłownym brzmieniu.
\end{enumerate}

\begin{lstlisting}[style=prompt,caption={Szablon promptu początkowego dla C (\file{C/README.md}).},label={lst:prompt-initial}]
You are a C programming expert. Your task is to implement the required functions in accordance with the provided header file.

[TASK-SPECIFIC REQUIREMENTS]

General Requirements:
- Memory must be managed dynamically (malloc/calloc). Functions responsible for deallocation (the `free_*` functions) must safely free all nested pointers and structures to avoid memory leaks.
- Return ONLY raw C code (no Markdown tags, no ```c blocks, no side text). Your code will be directly saved to a .c file and compiled. Do not provide main() or test functions.

[INSERT HEADER FILE CONTENT HERE]
\end{lstlisting}

\paragraph{Wskazówki dotyczące paradygmatu} Wymagania ogólne to miejsce, w którym szablony różnią się
między językami (\cref{tab:guidance}). W kilku słowach określają oczekiwania paradygmatu:
modelowanie obiektowe dla Javy i Smalltalka; czystość, niemodyfikowalność i
\code{Maybe}/\code{Option} dla Haskella i Scali; styl relacyjny i porażkę zamiast wyjątków dla
Prologu; dyscyplinę zarządzania pamięcią dla C i Pascala. Wskazówki są celowo krótkie. Chcemy
mierzyć, jakich idiomów model używa samodzielnie, a nie to, jak dobrze stosuje się do przewodnika
stylu. \Cref{app:prompts} zawiera wszystkie szablony i jeden kompletny wygenerowany prompt.

\begin{table}[htbp]
  \centering
  \caption{Specyficzne dla języka „General Requirements” promptów początkowych (pominięto
    instrukcje dotyczące formatu odpowiedzi, wspólne dla wszystkich języków).}
  \label{tab:guidance}
  \small
  \begin{tabularx}{\textwidth}{@{}lY@{}}
    \toprule
    Język & Wskazówki \\
    \midrule
    C         & Pamięć dynamiczna (\code{malloc}/\code{calloc}); funkcje \code{free_*} muszą zwalniać wszystkie zagnieżdżone struktury \\
    Pascal    & Idiomatyczny proceduralny Pascal; bezpieczne zarządzanie dynamicznie alokowanymi strukturami \\
    Java      & Idiomatyczna Java; projekt obiektowy, modelowanie dziedziny klasami; preferowane standardowe kolekcje \\
    Smalltalk & Idiomatyczny GNU Smalltalk; klasy i przekazywanie komunikatów; bez redefiniowania klas wbudowanych (\code{Point}, \code{Set}, \ldots); standardowe kolekcje \\
    Haskell   & Idiomatyczny, czysty Haskell; dane niemodyfikowalne, algebraiczne typy danych, funkcje totalne; \code{Maybe} zamiast null \\
    Scala     & Idiomatyczna Scala; niemodyfikowalne struktury danych i czyste funkcje; \code{Option} zamiast null \\
    Prolog    & Idiomatyczny SWI-Prolog; czyste predykaty i styl relacyjny; porażka zamiast wyjątków \\
    \bottomrule
  \end{tabularx}
\end{table}

\paragraph{Prompt naprawczy} Po nieudanej ewaluacji prompt naprawczy (\cref{lst:prompt-followup}) jest
dołączany do \emph{tej samej} konwersacji. Model widzi więc swoją wcześniejszą odpowiedź razem
z informacją zwrotną, podobnie jak programista czytający wyniki CI dla własnego commita. Informacją
zwrotną jest fragment logu i lista niezaliczonych testów (\cref{sec:repair-loop}). Prompt ponownie
wymaga kompletnego pliku i wyłącznie surowego kodu.

\begin{lstlisting}[style=prompt,caption={Szablon promptu naprawczego dla C.},label={lst:prompt-followup}]
Your previous C implementation failed the evaluation. Your task is to analyze the error and provide a corrected version of the code.

[ERROR LOGS / VALGRIND OUTPUT]

[FAILING TEST CASE (Code or name of the test case)]

Requirements for the fix:
1. Identify and resolve the core logical or memory issue causing the failure.
2. Ensure memory is managed safely.
3. Return ONLY the complete, raw, fixed C code for `[TARGET FILE NAME]`.
4. Do NOT include Markdown formatting, do NOT include ```c tags, and do NOT write any explanations. Just the code.
\end{lstlisting}

\paragraph{Zasady projektowe} Projekt promptów opiera się na kilku zasadach. (i) \emph{Testy
pozostają ukryte}, więc model nie może napisać kodu dopasowanego wyłącznie do testów. W rundzie
naprawy widzi nazwy testów i ich wyjście, ale nie ich kod. (ii) \emph{Specyfikacja jest kompletna}:
każde zachowanie sprawdzane przez test jest opisane w wymaganiach lub wynika z kontraktu.
(iii) Prompty są \emph{równoległe między językami}: ta sama struktura, to samo sformułowanie
z wyjątkiem języka i wskazówek dla niego, tak aby różnice w wynikach wynikały z języka, a nie
z promptu. (iv) \emph{Odpowiedź jest plikiem}: model zwraca kompletny plik docelowy, który jest
zapisywany na dysk bez zmian, z wyjątkiem usunięcia znaczników bloków kodu. (v) \emph{Generowanie
jest ścisłe}: szablon z niewypełnionym symbolem zastępczym powoduje błąd zamiast wysłania wadliwego
promptu.

\section{Generowanie i wycinanie kodu}
\label{sec:extraction}

Generator jest wywoływany przez strumieniowe Messages API~\cite{anthropic2026docs} z \emph{adaptacyjnym
myśleniem} i określonym poziomem \emph{wysiłku} (ang.~\emph{effort}; \cref{lst:generate}). Adaptacyjne
myślenie pozwala modelowi samodzielnie zdecydować, ile rozumować przed udzieleniem odpowiedzi.
Rozumowanie jest zwracane w osobnych blokach thinking, które przechowujemy w formie podsumowań do
późniejszej analizy. Obecne modele Claude nie przyjmują parametru \code{temperature}, więc odpowiedzi
różnią się między przebiegami. Z tego powodu system obsługuje wielokrotne \emph{próbki} na zadanie;
raportowany tu przebieg używa jednej próbki (\cref{sec:threats}). Cała odpowiedź, łącznie z blokami
thinking, jest dołączana do konwersacji bez zmian, aby rundy naprawy kontynuowały tok rozumowania
samego modelu.

\begin{lstlisting}[language=Python,caption={Żądanie generowania (\file{harness/llm.py}, w skrócie).},label={lst:generate}]
request = dict(
    model=model,                 # "claude-haiku-5-5"
    max_tokens=max_tokens,       # 32000 per response; thinking counts toward it
    thinking={"type": "adaptive", "display": "summarized"},
    output_config={"effort": effort},   # "medium"
    messages=list(messages),     # the whole conversation so far
)
response, request_id, duration = tracing.stream_message(client.messages.stream, request, trace_ctx)
code, violation = extract_code(response_text(response))  # text blocks only, never thinking
\end{lstlisting}

Prompty wymagają surowego kodu, ale odpowiedź i tak jest sprawdzana (\cref{lst:extract}). Jeśli
zawiera znaczniki bloków Markdown, wybierany jest największy blok. Jeśli brakuje znacznika
zamykającego, np.\ z powodu ucięcia odpowiedzi, wybierane jest wszystko po pierwszym znaczniku. Taka
odpowiedź jest oznaczana jako \emph{naruszenie formatu}. Naruszenie nie powoduje niezaliczenia
zadania, ale jest raportowane jako miara wykonywania instrukcji (\cref{sec:cost}). Gdy znaczników
nie ma wcale, tekst spoza kodu nie jest usuwany. Proza wpleciona w kod bez znaczników trafia więc do
kompilatora i powoduje niepowodzenie próby, co pokazują dwie odpowiedzi z naszego przebiegu
(\cref{sec:failure-analysis}).

\begin{lstlisting}[language=Python,caption={Wycinanie kodu i wykrywanie naruszeń formatu (\file{harness/llm.py}).},label={lst:extract}]
def extract_code(text: str) -> tuple[str, bool]:
    """Returns (code, format_violation). The prompts ask for raw code without Markdown fences."""
    stripped = text.strip()
    if not _FENCE_LINE.search(stripped):
        return (stripped + "\n" if stripped else ""), False
    blocks = _FENCED.findall(stripped)
    if blocks:
        return max(blocks, key=len).strip("\n") + "\n", True
    # Closing fence glued to the last code line, or missing (output cut off): take everything
    # between the first fence line and the last ``` (or the end), dropping prose before it.
    body = stripped[_FENCE_LINE.search(stripped).start():].split("\n", 1)
    body = body[1] if len(body) > 1 else ""
    end = body.rfind("```")
    if end >= 0:
        body = body[:end]
    return body.strip("\n") + "\n", True
\end{lstlisting}

\section{Ewaluacja i kategorie wyników}
\label{sec:outcomes}

Dla każdej próby katalog zadania jest kopiowany do katalogu roboczego ze znormalizowanymi końcami
linii i bez artefaktów budowania, a szkielet jest nadpisywany kodem modelu. Następnie
\file{run-tests.sh} uruchamia się w świeżym kontenerze z obrazem danego języka, z limitem czasu
300\,s. Wyjście jest przeszukiwane pod kątem linii \code{[PASS]} i \code{[FAIL]}; test, który nie
wypisał wyniku, liczy się jako niezaliczony. Każda próba otrzymuje dokładnie jedną kategorię wyniku
(\cref{tab:outcomes}). O pierwszych sześciu rozstrzygają znaczniki protokołu wyjścia, kod wyjścia
i limit czasu; ostatnie trzy oznaczają, że nie było kodu do oceny.

\begin{table}[htbp]
  \centering
  \caption{Kategorie wyniku próby. Zadanie jest \emph{rozwiązane}, gdy któraś próba osiąga
    \code{pass}.}
  \label{tab:outcomes}
  \small
  \begin{tabularx}{\textwidth}{@{}lY@{}}
    \toprule
    Kategoria & Znaczenie (reguła) \\
    \midrule
    \code{pass}          & wszystkie testy przeszły, a przebieg zakończył się kodem 0 \\
    \code{test_failures} & co najmniej jeden test nie przeszedł lub nie wypisał wyniku \\
    \code{memory_error}  & wszystkie testy przeszły, ale kod wyjścia jest niezerowy (Valgrind wykrył wycieki lub niepoprawne dostępy) \\
    \code{compile_error} & brak znacznika \code{=== START ===}: kod się nie skompilował lub nie załadował, żaden test nie został uruchomiony \\
    \code{runtime_crash} & brak znacznika \code{=== BENCHMARK RESULTS ===}: program testowy uległ awarii w trakcie testów \\
    \code{timeout}       & przebieg przekroczył limit czasu i kontener został zabity \\
    \code{truncated}     & odpowiedź osiągnęła limit tokenów wyjściowych, zanim kod był kompletny \\
    \code{refusal}       & model odmówił odpowiedzi \\
    \code{empty_output}  & odpowiedź nie zawierała kodu \\
    \bottomrule
  \end{tabularx}
\end{table}

\section{Pętla samonaprawy}
\label{sec:repair-loop}

Jeśli próba nie rozwiązała zadania, a pozostały rundy naprawy (co najwyżej trzy), system tworzy
\emph{fragment logu} dopasowany do kategorii wyniku i wysyła prompt naprawczy. Fragment zawiera
informacje, na które programista spojrzałby najpierw, ograniczone do 6\,000 znaków:
\begin{itemize}
  \item \code{compile_error}: początek wyjścia kompilatora lub programu ładującego, bez szumu
    budowania, np.\ linii \code{make clean};
  \item \code{test_failures}: ostrzeżenia kompilatora i diagnostykę ładowania, wszystkie linie
    \code{[PASS]}/\code{[FAIL]} wraz z komunikatami wypisanymi przez program testowy oraz linie
    z problemami zgłoszonymi przez Valgrind, jeśli takie są;
  \item \code{memory_error}: raport Valgrinda;
  \item \code{runtime_crash} i \code{timeout}: wyjaśnienie i końcówkę wyjścia.
\end{itemize}
Pętla (\cref{lst:loop}) kończy się przy pierwszym \code{pass}, przy odmowie lub po wyczerpaniu rund.
Każda próba jest zapisywana wraz z promptem, surową odpowiedzią, podsumowaniem rozumowania,
wyciętym kodem, logiem testów i wynikiem, a próbka jest podsumowywana w pliku \file{final.json}.

\begin{lstlisting}[language=Python,caption={Pętla generowanie--ewaluacja--naprawa (\file{harness/run.py}, w uproszczeniu).},label={lst:loop}]
messages, excerpt, failing = [], "", []
for attempt in range(cfg["max_repairs"] + 1):
    prompt = (prompts.render_initial(task) if attempt == 0
              else prompts.render_followup(task, excerpt, failing))
    messages.append({"role": "user", "content": prompt})
    gen = llm.generate(client, cfg["model"], cfg["effort"], messages, cfg["max_tokens"])
    messages.append({"role": "assistant", "content": gen.content})   # thinking included

    if gen.stop_reason == "refusal":    category = "refusal"
    elif not gen.code.strip():          category = "empty_output"
    elif gen.stop_reason == "max_tokens": category = "truncated"
    else:
        mount = evaluator.prepare_workdir(task, work, {task.target: gen.code})
        result, output = evaluator.run_tests(task, mount, timeout_s)
        category = result.category
        excerpt = evaluator.log_excerpt(result, output, timeout_s)
        failing = result.failing_names() if result.names else []
    save_attempt(...)
    if category in ("pass", "refusal"):
        break
\end{lstlisting}

\section{LLM jako sędzia idiomatyczności}
\label{sec:judge}

Testy mierzą poprawność; idiomatyczność wymaga czytelnika. Po zakończeniu generowania system wysyła
kod \emph{ostatniej} próby każdego zadania do drugiego modelu, \sonnet{} z wysokim poziomem
wysiłku, razem ze specyfikacją zadania, kontraktem i sekcją rubryki dla danego języka. Oceniany jest
tylko kod, który się skompilował.

\paragraph{Zaślepienie} Sędzia nie widzi wyników testów, liczby prób ani nazwy modelu, który napisał
kod. Prompt systemowy czyni go ekspertem recenzującym kod w danym języku i paradygmacie i poleca mu
oceniać styl, a nie poprawność (\cref{app:judge}).

\paragraph{Rubryka} Rubryka (\file{harness/rubrics.md}, przytoczona w \cref{app:judge}) definiuje
pięciostopniową skalę od 1 („walczy z językiem”) do 5 („tak napisałby to doświadczony praktyk”)
oraz cztery kryteria:
\begin{itemize}
  \item \emph{idiomatyczność paradygmatu}: dekompozycja, przepływ danych i struktury sterujące
    preferowane przez paradygmat;
  \item \emph{idiomatyczność języka}: wykorzystanie cech języka, jego biblioteki standardowej
    i konwencji;
  \item \emph{czytelność}: nazwy, rozmiar funkcji, struktura, przydatne komentarze;
  \item \emph{projekt}: dekompozycja, struktury danych, rozdzielenie odpowiedzialności, brak
    powtórzeń.
\end{itemize}
Sekcja dla każdego języka wymienia konkretne oczekiwania, np.\ sprawdzanie wyników alokacji
i funkcje pomocnicze \code{static} w C, protokoły kolekcji (\code{collect:}, \code{inject:into:})
w Smalltalku, funkcje totalne i \code{foldl'} w Haskellu, „zielone” odcięcia i akumulatory
w Prologu. Reguły ogólne zabraniają nagradzania długości, karania za wymagania wynikające z zadania
(struktury wynikowe w stylu C, atomy \code{none}, kolekcje Smalltalka indeksowane od 1) oraz
oceniania tego, co zostało dane w kontrakcie.

\paragraph{Ustrukturyzowana odpowiedź} Odpowiedź musi być zgodna ze schematem JSON
(\cref{lst:schema}). Kolejność właściwości ma znaczenie: model najpierw pisze \emph{uzasadnienie
dla każdego kryterium}, wskazując konkretny kod stojący za oceną i to, co dałoby jeden punkt więcej,
a dopiero potem oceny. Czytelnik może więc zweryfikować każdą ocenę względem kodu. Ocena zawiera
też listę antywzorców i krótkie uzasadnienie ogólne.

\begin{lstlisting}[language=JSON,caption={Schemat JSON oceny (w skrócie; \file{harness/judge.py}).},label={lst:schema}]
{ "type": "object",
  "properties": {
    "rationale": { "type": "object",            // written first, one entry per criterion
      "properties": { "paradigm_idiomaticity": {"type": "string"},
                      "language_idiomaticity": {"type": "string"},
                      "readability": {"type": "string"}, "design": {"type": "string"} } },
    "paradigm_idiomaticity": {"type": "integer", "enum": [1, 2, 3, 4, 5]},
    "language_idiomaticity": {"type": "integer", "enum": [1, 2, 3, 4, 5]},
    "readability":           {"type": "integer", "enum": [1, 2, 3, 4, 5]},
    "design":                {"type": "integer", "enum": [1, 2, 3, 4, 5]},
    "anti_patterns": {"type": "array", "items": {"type": "string"}},
    "justification": {"type": "string"} },
  "required": ["rationale", "paradigm_idiomaticity", "language_idiomaticity",
               "readability", "design", "anti_patterns", "justification"],
  "additionalProperties": false }
\end{lstlisting}

\paragraph{Kalibracja} Tryb \code{judge.py --calibrate} wielokrotnie ocenia ten sam kod (rozwiązania
wzorcowe i kilka rozwiązań z przebiegu, każde dwukrotnie), aby sprawdzić stabilność sędziego.
Wyniki znajdują się w \cref{sec:results-idiomaticity}.

\section{Powtarzalność, śledzenie i kontrola kosztów}
\label{sec:reproducibility}

\begin{itemize}
  \item \textbf{Wznawialne przebiegi.} Próbka z plikiem \file{final.json} jest pomijana przy
    ponownym uruchomieniu polecenia, a przerwana próbka jest wykonywana od nowa. Wyniki sędziego
    są pomijane w ten sam sposób. Pozwoliło to dodać zadania 7--10 do istniejącego przebiegu po ich
    napisaniu.
  \item \textbf{Ochrona konfiguracji.} Przebieg odmawia kontynuacji, jeśli model, poziom wysiłku,
    liczba rund naprawy lub limit tokenów różnią się od zapisanych w jego \file{config.json}.
    Zmiana ustawień wymaga nowego identyfikatora przebiegu.
  \item \textbf{Pochodzenie danych.} \file{config.json} rejestruje commit gita i to, czy drzewo
    robocze zawierało niezatwierdzone zmiany. Każde wywołanie API jest śledzone wraz z żądaniem,
    odpowiedzią, podsumowaniem rozumowania, identyfikatorem żądania, czasem trwania i zużyciem
    tokenów.
  \item \textbf{Kontrola kosztów.} Przed pierwszym płatnym żądaniem system wypisuje maksymalną
    liczbę wywołań i pesymistyczny szacunek kosztów oraz prosi o potwierdzenie.
  \item \textbf{Równoległość i izolacja.} Cztery zadania działają równolegle, każde we własnym
    kontenerze. Błędy krytyczne API (uwierzytelnianie, nieznany model, niepoprawne żądanie)
    przerywają przebieg; błędy przejściowe pozostawiają próbkę niedokończoną do następnego
    uruchomienia.
\end{itemize}

% =============================================================================
\chapter{Wyniki}
\label{ch:results}
% =============================================================================

Ten rozdział przedstawia wyniki przebiegu \code{haiku-v1}. Najpierw opisujemy konfigurację, a potem
kolejno odpowiadamy na pytania badawcze: poprawność (\RQ{1}, \cref{sec:results-correctness}),
rodzaje błędów i defekty benchmarku (\cref{sec:failure-analysis}), samonaprawa (\RQ{2},
\cref{sec:results-repair}), idiomatyczność (\RQ{3}, \cref{sec:results-idiomaticity}) oraz koszty
i wykonywanie instrukcji (\RQ{4}, \cref{sec:cost}).

\section{Konfiguracja eksperymentu}
\label{sec:setup}

\Cref{tab:setup} podsumowuje konfigurację. Każdą z 70 par język--zadanie uruchomiono raz (jedna
próbka). Zadania 1--6 uruchomiono 25~października~2026 na commicie \code{072ea49}. Zadania 7--10
dodano do tego samego przebiegu 13~listopada~2026, po ich scaleniu (commit \code{dc7f1d4}). Ochrona
konfiguracji (\cref{sec:reproducibility}) zapewniła identyczne ustawienia generatora. Wszystkie 70
rozwiązań końcowych oceniono z tą samą konfiguracją sędziego (schemat z uzasadnieniem przed
oceną).

\begin{table}[htbp]
  \centering
  \caption{Konfiguracja przebiegu \code{haiku-v1}.}
  \label{tab:setup}
  \small
  \begin{tabular}{@{}lll@{}}
    \toprule
    & Ustawienie & Wartość \\
    \midrule
    \multirow{5}{*}{Generator}
      & Model & \haiku{} (\code{claude-haiku-5-5}) \\
      & Rozumowanie & adaptacyjne myślenie, wysiłek \code{medium}, podsumowania \\
      & Maks.\ tokenów wyjściowych & 32\,000 na odpowiedź (łącznie z rozumowaniem) \\
      & Rundy naprawy & do 3 (co najwyżej 4 próby na zadanie) \\
      & Próbki na zadanie & 1 \\
    \midrule
    \multirow{4}{*}{Sędzia}
      & Model & \sonnet{} (\code{claude-sonnet-5-5}) \\
      & Rozumowanie & adaptacyjne myślenie, wysiłek \code{high} \\
      & Oceniana próba & ostatnia próba każdego zadania (jeśli się skompilowała) \\
      & Odpowiedź & schemat JSON, uzasadnienie przed ocenami (\cref{lst:schema}) \\
    \midrule
    \multirow{3}{*}{Ewaluacja}
      & Limit czasu & 300\,s na przebieg testów, 4 zadania równolegle \\
      & Skala & 7 języków $\times$ 10 zadań $=$ 70 zadań, 1\,400 testów \\
      & Wywołania API & 98 wywołań generujących, 70 wywołań sędziego \\
    \bottomrule
  \end{tabular}
\end{table}

% -----------------------------------------------------------------------------
\section{RQ1: Poprawność funkcjonalna}
\label{sec:results-correctness}

\Cref{tab:matrix} pokazuje wynik każdego zadania, a \cref{tab:correctness} agreguje go według
języka i paradygmatu. Łącznie model rozwiązał 56 z 70 zadań (80,0\,\%) pierwszą odpowiedzią
i 65 (92,9\,\%) po samonaprawie. W ujęciu testów zaliczył 1\,229 z 1\,400 testów (87,8\,\%) za
pierwszym podejściem i 1\,375 (98,2\,\%) na końcu. W odróżnieniu od pilotażu z HumanEval
(\cref{sec:humaneval}) benchmark pozostawia więc miejsce na obserwację różnic.

\begin{table}[htbp]
  \centering
  \caption{Wyniki dla poszczególnych języków i zadań. Komórka pokazuje liczbę zaliczonych testów
    (z 20) w pierwszej próbie oraz, jeśli jest inna, w próbie ostatniej. \cmark{} oznacza zadanie
    rozwiązane, a indeks górny rundę naprawy, w której je rozwiązano; zacienione komórki z \xmark{}
    nie zostały rozwiązane po trzech rundach naprawy.}
  \label{tab:matrix}
  \footnotesize
  \setlength{\tabcolsep}{3.2pt}
  \begin{tabular}{@{}l*{10}{c}c@{}}
    \toprule
    Język & Z1 & Z2 & Z3 & Z4 & Z5 & Z6 & Z7 & Z8 & Z9 & Z10 & Rozw. \\
    \midrule
    \tabinput{tables/matrix}
    \bottomrule
  \end{tabular}
\end{table}

\begin{table}[htbp]
  \centering
  \caption{Poprawność według języka i paradygmatu: odsetek zaliczonych testów i liczba
    rozwiązanych zadań w pierwszej próbie i po samonaprawie (końcowo).}
  \label{tab:correctness}
  \small
  \begin{tabular}{@{}llrrrr@{}}
    \toprule
    & & \multicolumn{2}{c}{Zaliczone testy (\%)} & \multicolumn{2}{c}{Rozwiązane zadania} \\
    \cmidrule(lr){3-4}\cmidrule(lr){5-6}
    Język & Paradygmat & pierwsza & końcowo & pierwsza & końcowo \\
    \midrule
    C         & imperatywny & 100,0 & 100,0 & 10/10 & 10/10 \\
    Pascal    & imperatywny &  90,0 &  98,0 &  9/10 &  9/10 \\
    Java      & obiektowy   &  70,5 & 100,0 &  6/10 & 10/10 \\
    Smalltalk & obiektowy   &  85,5 &  96,0 &  8/10 &  9/10 \\
    Haskell   & funkcyjny   &  78,0 &  99,5 &  7/10 &  9/10 \\
    Scala     & funkcyjny   &  99,0 & 100,0 &  9/10 & 10/10 \\
    Prolog    & logiczny    &  91,5 &  94,0 &  7/10 &  8/10 \\
    \midrule
    \multicolumn{2}{@{}l}{Imperatywny (C, Pascal)}        & 95,0 & 99,0 & 19/20 & 19/20 \\
    \multicolumn{2}{@{}l}{Obiektowy (Java, Smalltalk)}    & 78,0 & 98,0 & 14/20 & 19/20 \\
    \multicolumn{2}{@{}l}{Funkcyjny (Haskell, Scala)}     & 88,5 & 99,8 & 16/20 & 19/20 \\
    \multicolumn{2}{@{}l}{Logiczny (Prolog)}              & 91,5 & 94,0 &  7/10 &  8/10 \\
    \midrule
    \multicolumn{2}{@{}l}{Łącznie}                        & 87,8 & 98,2 & 56/70 & 65/70 \\
    \bottomrule
  \end{tabular}
\end{table}

\begin{figure}[htbp]
  \centering
  % Share of tests passed at the first and at the final attempt, per language (dumbbell plot).
% Data: results/haiku-v1/summary.csv (first_n_pass, final_n_pass; 200 tests per language).
\begin{tikzpicture}
\begin{axis}[
  reportchart,
  width=0.86\textwidth, height=6.4cm,
  xmin=65, xmax=102, ymin=0.4, ymax=7.6,
  xtick={70,75,80,85,90,95,100},
  xlabel={Zaliczone testy (\%)},
  ytick={1,...,7},
  yticklabels={Prolog (log.), Scala (funk.), Haskell (funk.), Smalltalk (ob.), Java (ob.), Pascal (imp.), C (imp.)},
  xmajorgrids=true, ymajorgrids=false,
  axis x line*=bottom, axis y line*=left,
  y axis line style={draw=none}, ytick style={draw=none},
  legend style={at={(0.0,1.02)}, anchor=south west, legend columns=2, /tikz/every even column/.append style={column sep=10pt}},
  clip=false,
]
  % connectors (first -> final)
  \foreach \y/\a/\b in {7/100/100, 6/90.0/98.0, 5/70.5/100, 4/85.5/96.0, 3/78.0/99.5, 2/99.0/100, 1/91.5/94.0}{
    \edef\temp{\noexpand\draw[draw=neutral, line width=2pt] (axis cs:\a,\y) -- (axis cs:\b,\y);}
    \temp
  }
  \addplot[only marks, mark=*, mark size=3.4pt, mark options={fill=seq250, draw=white, line width=0.8pt}]
    coordinates {(100,7) (90.0,6) (70.5,5) (85.5,4) (78.0,3) (99.0,2) (91.5,1)};
  \addlegendentry{pierwsza próba}
  \addplot[only marks, mark=*, mark size=3.4pt, mark options={fill=seq600, draw=white, line width=0.8pt}]
    coordinates {(100,7) (98.0,6) (100,5) (96.0,4) (99.5,3) (100,2) (94.0,1)};
  \addlegendentry{po samonaprawie (końcowo)}
  % selective direct labels: the two largest repair gains
  \node[font=\scriptsize, text=inksec, anchor=south] at (axis cs:70.5,5.18) {70,5};
  \node[font=\scriptsize, text=inksec, anchor=south] at (axis cs:78.0,3.18) {78,0};
\end{axis}
\end{tikzpicture}

  \caption{Odsetek z 200 testów na język zaliczonych w pierwszej próbie (jasne punkty) i po
    samonaprawie (ciemne). Dla C oba punkty pokrywają się na poziomie 100\,\%.}
  \label{fig:passrate}
\end{figure}

\paragraph{Języki} C jest jedynym językiem, w którym każde zadanie przeszło za pierwszym podejściem,
łącznie ze wszystkimi kontrolami Valgrinda. Dotyczy to również wyniku bez 17 testów przechodzących
na szkielecie (183/183). Następna jest Scala z 9 zadaniami i 99\,\% testów. Java ma najsłabszą
pierwszą próbę: 6 zadań i 70,5\,\% testów (\cref{fig:passrate}). Jej błędy są jednak przeważnie
tanie do naprawienia: dwa z nich to błędy kompilacji spowodowane wymyśloną przez model deklaracją
\code{package} (\cref{sec:failure-analysis}), a wszystkie cztery nierozwiązane za pierwszym razem
zadania w Javie zostały naprawione. Prolog kończy z najniższym wynikiem: 8 rozwiązanych zadań
i 94,0\,\% testów, ale oba jego nierozwiązane zadania są dotknięte defektami w zestawach testów dla
Prologu (\cref{sec:benchmark-defects}).

\paragraph{Paradygmaty} Za pierwszym podejściem kolejność paradygmatów to: imperatywny (19/20 zadań),
funkcyjny (16/20), obiektowy (14/20) i logiczny (7/10). Po samonaprawie paradygmaty imperatywny,
obiektowy i funkcyjny osiągają po 19/20, a logiczny 8/10. Różnice \emph{wewnątrz} paradygmatu są
równie duże jak różnice między paradygmatami. Java rozwiązała za pierwszym podejściem dwa zadania
mniej niż Smalltalk, a Haskell dwa mniej niż Scala. Mniej popularny język z pary nie jest też
konsekwentnie gorszy: Smalltalk wypadł lepiej niż Java, a Pascal niemal dorównał C.

\paragraph{Ostrożność statystyczna} Przy jednej próbce dziesięciu zadań na język żadna z tych różnic
nie jest istotna statystycznie. Największa różnica w pierwszej próbie między dwoma językami, C
(10/10) wobec Javy (6/10), daje $p = 0,087$ w dwustronnym dokładnym teście Fishera. Paradygmat
imperatywny wobec obiektowego (19/20 wobec 14/20) daje $p = 0,091$, a imperatywny wobec logicznego
(19/20 wobec 7/10) $p = 0,095$. Wyniki należy traktować jako przesłanki uzasadniające powtórzone
próbki, a nie jako ustalone efekty (\cref{sec:threats}).

\paragraph{Zadania} Zadanie~1 (logika przestrzenna) okazało się najtrudniejsze: za pierwszym
podejściem rozwiązały je tylko 3 z 7 języków, głównie z powodu danych zdegenerowanych (pojedyncze
punkty, wielokąty puste i płaskie). Następne było zadanie~9 (B-drzewo, 4 z 7), ponieważ dokładny
kształt drzewa nie zostawia miejsca na warianty algorytmu. Zadania~6 (Sudoku) i~8 (TSP) zostały
rozwiązane za pierwszym podejściem we wszystkich siedmiu językach, z zaliczeniem wszystkich 140
testów. Zadanie~2 (Dijkstra) jest jedynym, którego nie rozwiązały dwa języki.

% -----------------------------------------------------------------------------
\section{Analiza błędów}
\label{sec:failure-analysis}

\Cref{tab:outcomes-lang} klasyfikuje wszystkie 98 prób. Występują tylko trzy z dziewięciu kategorii:
nie było żadnego błędu pamięci, awarii, przekroczenia limitu czasu, uciętej odpowiedzi, odmowy ani
pustej odpowiedzi. Wszystkie 98 odpowiedzi zakończyło się normalnie (powód zatrzymania
\code{end_turn}); najdłuższa wykorzystała 23\,034 z 32\,000 dostępnych tokenów wyjściowych.
Najwięcej niezaliczonych testów przypada na Prolog (9 prób), a najwięcej błędów kompilacji na
Javę (4).

\begin{table}[htbp]
  \centering
  \caption{Kategorie wyników wszystkich prób, łącznie z rundami naprawy. Pominięto kategorie, które
    nie wystąpiły (błąd pamięci, awaria w czasie wykonania, przekroczenie czasu, ucięcie, odmowa,
    pusta odpowiedź).}
  \label{tab:outcomes-lang}
  \small
  \begin{tabular}{@{}lrrrr@{}}
    \toprule
    Język & \code{pass} & \code{test_failures} & \code{compile_error} & Próby \\
    \midrule
    C         & 10 & 0 & 0 & 10 \\
    Pascal    &  9 & 3 & 1 & 13 \\
    Java      & 10 & 2 & 4 & 16 \\
    Smalltalk &  9 & 4 & 2 & 15 \\
    Haskell   &  9 & 4 & 2 & 15 \\
    Scala     & 10 & 1 & 1 & 12 \\
    Prolog    &  8 & 9 & 0 & 17 \\
    \midrule
    Łącznie   & 65 & 23 & 10 & 98 \\
    \bottomrule
  \end{tabular}
\end{table}

\paragraph{Błędy kompilacji} Przeanalizowaliśmy wszystkie dziesięć błędów kompilacji
(\cref{tab:compile-errors}). \emph{Połowa z nich w ogóle nie jest błędami programistycznymi}:
odpowiedź zawierała tekst niebędący kodem. W czterech rundach naprawy (Java, zadania 1 i 9; Scala,
zadanie 1; Smalltalk, zadanie 9) model zaczął odpowiedź od zdania wyjaśniającego błąd
(\cref{lst:leak}), mimo że prompt naprawczy mówi „do NOT write any explanations”. W jednej
pierwszej próbie (Haskell, zadanie 4) kod kończył się zabłąkanym znacznikiem HTML \code{</br>}.
Ponieważ te odpowiedzi nie miały znaczników bloków Markdown, mechanizm wycinania pozostawił je bez
zmian, a kompilator je odrzucił (\cref{sec:extraction}). Pozostałe pięć błędów kompilacji to
rzeczywiste pomyłki. Są to: deklaracja \code{package} niezgodna z domyślnym pakietem kontraktu
(Java, zadania 9 i 10), lokalna zmienna \code{L} pisana wielką literą, którą Haskell interpretuje
jako konstruktor danych (zadanie 3), wywołanie \code{IsInfinity} bez modułu \code{Math} (Pascal,
zadanie 5) oraz rozwiązanie w Smalltalku, którego klasy się nie załadowały (zadanie 2).

\begin{table}[htbp]
  \centering
  \caption{Przyczyny dziesięciu błędów kompilacji.}
  \label{tab:compile-errors}
  \small
  \begin{tabularx}{\textwidth}{@{}llcY@{}}
    \toprule
    Język & Zadanie & Próba & Przyczyna \\
    \midrule
    Java      & 1  & 1 & odpowiedź zaczyna się od wyjaśnienia prozą \\
    Java      & 9  & 1 & odpowiedź zaczyna się od wyjaśnienia prozą \\
    Scala     & 1  & 1 & odpowiedź zaczyna się od wyjaśnienia prozą \\
    Smalltalk & 9  & 1 & odpowiedź zaczyna się od wyjaśnienia prozą \\
    Haskell   & 4  & 0 & zabłąkany znacznik \code{</br>} po ostatniej linii kodu \\
    \midrule
    Java      & 9  & 0 & wymyślona deklaracja \code{package com.example.btree;} ukrywa klasę kontraktu \\
    Java      & 10 & 0 & wymyślona deklaracja \code{package com.example;} ukrywa klasę kontraktu \\
    Haskell   & 3  & 0 & lokalne wiązanie o nazwie \code{L}, interpretowane jako konstruktor danych \\
    Pascal    & 5  & 0 & \code{IsInfinity} użyte bez modułu \code{Math} \\
    Smalltalk & 2  & 0 & klasy nie załadowały się do obrazu (\texttt{nil doesNotUnderstand: \#new}) \\
    \bottomrule
  \end{tabularx}
\end{table}

\begin{lstlisting}[style=prompt,caption={Początek dwóch odpowiedzi naprawczych, które zostały zapisane jako pliki źródłowe (Java, zadanie 9; Smalltalk, zadanie 9; linie skrócone).},label={lst:leak}]
package com.example.btree is the root cause. BTree.java and Tests.java are in the default package, so BTreeImpl in a named package cannot see BTree and its inherited nested type Node. ...

The failure is in the deletion path. Every test that removes a key fails, while insertion and the absent-key test pass. The likely cause is `OrderedCollection>>...
\end{lstlisting}

Jest to ten sam problem, który uczynił modele lokalne bezużytecznymi (\cref{sec:local-llms}), ale
o znacznie mniejszej częstości: 4 z 28 odpowiedzi naprawczych (14\,\%) i 1 z 70 pierwszych
odpowiedzi. Prompt naprawczy wydaje się skłaniać do wyjaśnień, ponieważ model właśnie usłyszał, że
jego kod zawiódł. Te wycieki prozy zawyżają też liczbę błędów kompilacji w Javie, Scali i Smalltalku
i kosztują każdy z tych języków jedną rundę naprawy.

\paragraph{Zadania nierozwiązane} Pięć zadań pozostało nierozwiązanych:
\begin{itemize}
  \item \emph{Pascal, zadanie 5 (16/20).} Po dodaniu brakującego modułu cztery testy nadal nie
    przechodziły: zagnieżdżone operatory jednoargumentowe, jednoargumentowy minus a potęgowanie
    ($-2^2 = -(2^2)$), funkcje wbudowane i błędy dziedziny. Są to reguły priorytetów i semantyki
    błędów ze specyfikacji.
  \item \emph{Smalltalk, zadanie 2 (12/20).} Odtwarzanie ścieżki przechodzi, ale zapytania
    o odległości nie. Program testowy w Smalltalku przechwytuje każdy wyjątek i wypisuje tylko
    \code{[FAIL]}, więc prompty naprawcze nie zawierały żadnej diagnostyki poza nazwami testów.
    Liczba zaliczonych testów wzrosła z 3 do 12, po czym przestała się zmieniać.
  \item \emph{Haskell, zadanie 3 (19/20).} Jeden test, obsługa brakujących argumentów
    (\code{NULL}), nadal nie przechodził po trzech rundach.
  \item \emph{Prolog, zadania 1 i 2 (11/20 i 17/20).} Oba są dotknięte defektami w zestawach testów
    dla Prologu, opisanymi poniżej.
\end{itemize}

\subsection{Defekty w zestawach testów dla Prologu}
\label{sec:benchmark-defects}

Analiza dwóch nierozwiązanych zadań w Prologu wykazała, że błędy są, przynajmniej częściowo,
spowodowane przez benchmark, a nie przez model.

\paragraph{Zadanie 2: błąd o jeden w funkcji pomocniczej testów} Testy obciążeniowe A10, B05 i B10
budują łańcuch o 1\,000 wierzchołków za pomocą predykatu pomocniczego \code{create_linear_chain/3}.
Wersja dla Prologu iteruje po wszystkich wierzchołkach $0,\dots,N{-}1$ i dodaje krawędź od każdego
z nich do następnika, w tym krawędź $999 \to 1000$, która wskazuje poza graf o $N = 1000$
wierzchołkach. Funkcja pomocnicza w C kończy o jeden wierzchołek wcześniej (\cref{lst:defect}).
Specyfikacja dla Prologu nie mówi, co \code{add_edge/5} ma zrobić z wierzchołkiem spoza zakresu.
Implementacja modelu go odrzuca, co jest rozsądnym wyborem, więc graf testowy nigdy nie zostaje
zbudowany i wszystkie trzy testy nie przechodzą niezależnie od implementacji algorytmu Dijkstry.
Prompty naprawcze pokazywały tylko nazwy testów, więc model nie mógł odkryć przyczyny. Liczba
zaliczonych testów pozostała równa 17 we wszystkich czterech próbach. Pozostałe 17 testów przeszło
za pierwszym podejściem; czy rozwiązanie przechodzi też testy obciążeniowe, będzie wiadomo dopiero
po poprawieniu funkcji pomocniczej.

\begin{lstlisting}[language=Prolog,caption={Predykat pomocniczy budujący łańcuch w testach dla Prologu (u góry) dodaje krawędź spoza zakresu $N{-}1 \to N$; funkcja pomocnicza w C (u dołu) tego nie robi.},label={lst:defect}]
create_linear_chain(NumVertices, Weight, Graph) :-
    create_graph(NumVertices, G0),
    Max is NumVertices - 1,
    numlist(0, Max, Indices),                     % 0 .. N-1
    foldl(add_chain_edge(Weight), Indices, G0, Graph).
add_chain_edge(Weight, I, GIn, GOut) :-
    J is I + 1,                                   % for I = N-1: J = N (out of range)
    add_edge(GIn, I, J, Weight, GOut).

/* C: for (int i = 0; i < num_vertices - 1; i++) graph_add_edge(g, i, i + 1, weight); */
\end{lstlisting}

\paragraph{Zadanie 1: niedospecyfikowany kontrakt} W odróżnieniu od późniejszych kontraktów dla
Prologu (\cref{lst:contract-pl}) szkielet zadania 1 nie dokumentuje ani nazw argumentów, ani ich
trybów: \code{check_dynamic_collision(_, _, _)}. Jego komentarz zapisuje też strukturę termów
z nazwami pisanymi wielką literą (\code{Polygon(...)}), które różnią się od termów używanych
w testach (\code{polygon(...)}). Testy wywołują \code{check_dynamic_collision(Obstacle, Robot,
Velocity)}, zgodnie z nagłówkiem w C, podczas gdy specyfikacja w języku naturalnym wymienia robota
jako pierwszego. Model wybrał kolejność \code{(Robot, Obstacle, Velocity)}, która zamienia miejscami
wielokąt ruchomy i nieruchomy, przez co wszystkie osiem nietrywialnych testów kolizji (B01--B08) nie
przechodzi. Test A08 oczekuje ponadto, że argumenty \code{none} dadzą wynik \code{none}, o czym
specyfikacja dla Prologu nie wspomina. W C, Javie, Haskellu i Scali nazwy lub typy parametrów
rozstrzygają kolejność wywołania i żaden inny język nie oblał testów kolizji.

\begin{findingbox}[Spostrzeżenie.]
Oba nierozwiązane zadania w Prologu pochodzą z wczesnych, ręcznie napisanych zestawów testów;
generowane zadania 6--10 nie mają takich problemów. Surowe liczby w tym rozdziale podajemy tak, jak
zostały zmierzone, ale wyniki Prologu dla zadań 1 i 2 zaniżają możliwości modelu. Zestawy te należy
poprawić przed kolejnymi przebiegami (\cref{ch:conclusions}).
\end{findingbox}

% -----------------------------------------------------------------------------
\section{RQ2: Samonaprawa}
\label{sec:results-repair}

\Cref{tab:repair} i \cref{fig:repair} pokazują, kiedy rozwiązano każde zadanie. Spośród 14 zadań
nierozwiązanych za pierwszym podejściem naprawiono 9 (64\,\%): 5 w pierwszej rundzie i 4 w drugiej.
Trzecia runda nigdy nie pomogła. Naprawione zadanie wymagało średnio 1,4 rundy.

\begin{table}[htbp]
  \centering
  \caption{Samonaprawa według języka. „Naprawione” to zadania, które nie przeszły w próbie 0, ale
    zostały rozwiązane później; trajektoria podaje liczbę zaliczonych testów w kolejnych próbach
    (C = błąd kompilacji, P = zaliczone).}
  \label{tab:repair}
  \small
  \begin{tabularx}{\textwidth}{@{}lccccY@{}}
    \toprule
    Język & W próbie 0 & Naprawione & Nierozw. & Śr.\ runda & Trajektorie nieudanych zadań \\
    \midrule
    C         & 10 & --       & 0 & --  & -- \\
    Pascal    &  9 & 0/1      & 1 & --  & Z5: 0C, 16, 16, 16 \\
    Java      &  6 & 4/4      & 0 & 1,5 & Z1: 19, 0C, 20P; Z7: 2, 20P; Z9: 0C, 0C, 20P; Z10: 0C, 20P \\
    Smalltalk &  8 & 1/2      & 1 & 2,0 & Z2: 0C, 3, 12, 12; Z9: 11, 0C, 20P \\
    Haskell   &  7 & 2/3      & 1 & 1,0 & Z1: 16, 20P; Z3: 0C, 19, 19, 19; Z4: 0C, 20P \\
    Scala     &  9 & 1/1      & 0 & 2,0 & Z1: 18, 0C, 20P \\
    Prolog    &  7 & 1/3      & 2 & 1,0 & Z1: 10, 11, 11, 11; Z2: 17, 17, 17, 17; Z9: 16, 20P \\
    \midrule
    Łącznie   & 56 & 9/14     & 5 & 1,4 & \\
    \bottomrule
  \end{tabularx}
\end{table}

\begin{figure}[htbp]
  \centering
  % Attempt at which each task was solved, per language (stacked bars, 10 tasks per language).
% Data: results/haiku-v1/report.md, section 4 (self-repair).
\begin{tikzpicture}
\begin{axis}[
  reportchart,
  xbar stacked,
  width=0.86\textwidth, height=6.2cm,
  bar width=10pt,
  xmin=0, xmax=10, xtick={0,2,...,10},
  xlabel={Zadania (z 10)},
  ytick={1,...,7}, ymin=0.4, ymax=7.6,
  yticklabels={Prolog, Scala, Haskell, Smalltalk, Java, Pascal, C},
  xmajorgrids=true,
  axis x line*=bottom, axis y line*=left,
  y axis line style={draw=none}, ytick style={draw=none},
  legend style={at={(0.0,1.02)}, anchor=south west, legend columns=4,
                /tikz/every even column/.append style={column sep=8pt}},
  every axis plot/.append style={draw=white, line width=0.8pt},
  legend image code/.code={\draw[#1, draw=none] (0cm,-0.1cm) rectangle (0.3cm,0.12cm);},
]
  \addplot[fill=seq600] coordinates {(10,7) (9,6) (6,5) (8,4) (7,3) (9,2) (7,1)};
  \addlegendentry{próba 0 (pierwsza odpowiedź)}
  \addplot[fill=seq400] coordinates {(0,7) (0,6) (2,5) (0,4) (2,3) (0,2) (1,1)};
  \addlegendentry{runda naprawy 1}
  \addplot[fill=seq250] coordinates {(0,7) (0,6) (2,5) (1,4) (0,3) (1,2) (0,1)};
  \addlegendentry{runda naprawy 2}
  \addplot[fill=neutral] coordinates {(0,7) (1,6) (0,5) (1,4) (1,3) (0,2) (2,1)};
  \addlegendentry{nierozwiązane}
\end{axis}
\end{tikzpicture}

  \caption{Próba, w której rozwiązano każde z dziesięciu zadań w danym języku.}
  \label{fig:repair}
\end{figure}

\paragraph{Co udaje się naprawić} Naprawa działa najlepiej, gdy informacja zwrotna wskazuje
bezpośrednio na defekt. Poza Prologiem wszystkie pięć zadań, które w pierwszej próbie oblały testy,
zostało naprawionych w ciągu jednej lub dwóch rund na podstawie linii \code{[FAIL]}
i komunikatów programu testowego; w Prologu udało się to w jednym przypadku na trzy. Z sześciu błędów
kompilacji w pierwszej próbie pięć skompilowało się po jednej rundzie, ponieważ komunikat
kompilatora nazywa problem; szósty (Java, zadanie 9) skompilował się po drugiej rundzie, gdyż
pierwsza odpowiedź naprawcza zaczynała się od prozy. Tylko trzy z sześciu przeszły następnie
wszystkie testy. Pozostałe błędy miały charakter semantyczny (reguły priorytetów w Pascalu) lub
były niewidoczne w informacji zwrotnej.

\paragraph{Kiedy naprawa utyka} We wszystkich pięciu nierozwiązanych zadaniach liczba zaliczonych
testów przestała się zmieniać po co najwyżej dwóch rundach (\cref{tab:repair}). W każdym przypadku
informacja zwrotna niosła niewiele treści. Defektów w Prologu nie da się dostrzec z nazw testów,
program testowy w Smalltalku połyka wyjątki, a jedyny pozostały błąd w Haskellu dotyczy konwencji
(brakujących argumentów), której log nie wyjaśnia. Jest to zgodne z ustaleniami
\citet{olausson2024selfrepair}, którzy wskazali jakość informacji zwrotnej jako wąskie gardło
samonaprawy. Dla twórców benchmarków oznacza to, że programy testowe powinny wypisywać wartość
oczekiwaną i rzeczywistą dla niespełnionej asercji.

\paragraph{Regresje} Trzy naprawy przeszły przez błąd kompilacji, zanim się powiodły (Java,
zadanie 1; Scala, zadanie 1; Smalltalk, zadanie 9). W każdej z nich pierwsza odpowiedź naprawcza
zaczynała się od prozy (\cref{tab:compile-errors}), a druga ją usunęła. Bez tych błędów formatu
zadania te zostałyby prawdopodobnie naprawione o rundę wcześniej, a średnia runda naprawy byłaby
niższa.

\paragraph{Paradygmaty} Odsetki naprawy dla paradygmatów wynoszą 0/1 (imperatywny), 5/6 (obiektowy),
3/4 (funkcyjny) i 1/3 (logiczny). Liczby są zbyt małe, by je porównywać. Niski wynik paradygmatu
logicznego tłumaczą dwa defekty w zestawach testów, a wynik obiektowy jest zdominowany przez tanie
poprawki błędów kompilacji w Javie.

% -----------------------------------------------------------------------------
\section{RQ3: Idiomatyczność}
\label{sec:results-idiomaticity}

\Cref{tab:judge} zestawia średnie oceny sędziego według języka i paradygmatu, a \cref{fig:judge}
porównuje paradygmaty.

\begin{table}[htbp]
  \centering
  \caption{Średnie oceny idiomatyczności rozwiązań końcowych (1 = słabo, 5 = doskonale), wystawione
    na ślepo przez \sonnet{}. $n$ = liczba ocenionych rozwiązań; AW = średnia liczba antywzorców
    wskazanych na rozwiązanie. Najlepsza wartość w kolumnie spośród języków jest pogrubiona.}
  \label{tab:judge}
  \small
  \begin{tabular}{@{}lrccccc r@{}}
    \toprule
    & $n$ & Paradygmat & Język & Czytelność & Projekt & Średnia & AW \\
    \midrule
    C         & 10 & 4,10 & 3,50 & 3,60 & 3,40 & 3,65 & 4,6 \\
    Pascal    & 10 & 4,00 & 3,40 & 3,50 & 3,50 & 3,60 & 4,6 \\
    Java      & 10 & 3,50 & 3,80 & 3,70 & 3,60 & 3,65 & 4,0 \\
    Smalltalk & 10 & 3,00 & 3,30 & 3,50 & 3,20 & 3,25 & 4,8 \\
    Haskell   & 10 & 4,20 & 3,70 & 3,70 & \textbf{3,70} & 3,83 & 3,9 \\
    Scala     & 10 & \textbf{4,30} & \textbf{4,00} & \textbf{3,90} & \textbf{3,70} & \textbf{3,98} & \textbf{3,1} \\
    Prolog    & 10 & 3,40 & 3,40 & 3,50 & 3,30 & 3,40 & 5,0 \\
    \midrule
    Imperatywny & 20 & 4,05 & 3,45 & 3,55 & 3,45 & 3,62 & 4,6 \\
    Obiektowy   & 20 & 3,25 & 3,55 & 3,60 & 3,40 & 3,45 & 4,4 \\
    Funkcyjny   & 20 & 4,25 & 3,85 & 3,80 & 3,70 & 3,90 & 3,5 \\
    Logiczny    & 10 & 3,40 & 3,40 & 3,50 & 3,30 & 3,40 & 5,0 \\
    \bottomrule
  \end{tabular}
\end{table}

\begin{figure}[htbp]
  \centering
  % Mean judge scores per paradigm and criterion (1-5 scale; 20 solutions per paradigm, 10 for logic).
% Data: results/haiku-v1/report.md, section 5. Colours: categorical slots 1-4 in fixed order.
\begin{tikzpicture}
\begin{axis}[
  reportchart,
  ybar=1.5pt,
  width=0.9\textwidth, height=6.2cm,
  bar width=11pt,
  ymin=1, ymax=5, ytick={1,2,3,4,5},
  ylabel={Średnia ocena (1--5)},
  symbolic x coords={P, L, R, D},
  xtick=data,
  xticklabels={Idiomy paradygmatu, Idiomy języka, Czytelność, Projekt},
  enlarge x limits=0.16,
  ymajorgrids=true,
  axis x line*=bottom, axis y line*=left,
  x axis line style={draw=inkmuted},
  legend style={at={(0.0,1.02)}, anchor=south west, legend columns=4,
                /tikz/every even column/.append style={column sep=8pt}},
  every axis plot/.append style={draw=white, line width=0.6pt},
  legend image code/.code={\draw[#1, draw=none] (0cm,-0.1cm) rectangle (0.3cm,0.12cm);},
]
  \addplot[fill=cat1] coordinates {(P,4.05) (L,3.45) (R,3.55) (D,3.45)};
  \addlegendentry{imperatywny}
  \addplot[fill=cat2] coordinates {(P,3.25) (L,3.55) (R,3.60) (D,3.40)};
  \addlegendentry{obiektowy}
  \addplot[fill=cat3] coordinates {(P,4.25) (L,3.85) (R,3.80) (D,3.70)};
  \addlegendentry{funkcyjny}
  \addplot[fill=cat4] coordinates {(P,3.40) (L,3.40) (R,3.50) (D,3.30)};
  \addlegendentry{logiczny}
\end{axis}
\end{tikzpicture}

  \caption{Średnie oceny sędziego według paradygmatu i kryterium. Oś zaczyna się od 1, najniższej
    oceny w rubryce.}
  \label{fig:judge}
\end{figure}

\paragraph{Kod funkcyjny jest najbardziej idiomatyczny} Paradygmat funkcyjny prowadzi we wszystkich
kryteriach (średnia 3,90), przy czym Scala jest najlepszym językiem (3,98) i ma najmniej
antywzorców na rozwiązanie (3,1). Dalej plasuje się paradygmat imperatywny (3,62), następnie
obiektowy (3,45) i logiczny (3,40). Smalltalk ma najniższą ocenę spośród wszystkich języków (3,25).

\paragraph{Idiomatyczność paradygmatu najsilniej różnicuje paradygmaty} Kryterium mierzące sam
paradygmat ma największy rozrzut, od 3,00 (Smalltalk) do 4,30 (Scala). Kod imperatywny (4,05)
i funkcyjny (4,25) uzyskuje wysokie oceny, obiektowy (3,25) i logiczny (3,40) niskie. Smalltalk
otrzymał za paradygmat dokładnie 3 we wszystkich dziesięciu zadaniach (\cref{tab:paradigm-scores}
w \cref{app:results}). Antywzorce wskazane przez sędziego (\cref{tab:antipatterns}) pokazują
dlaczego. W językach obiektowych model pisze kod proceduralny w klasach: statyczne metody pomocnicze
zamiast obiektów dziedzinowych w Javie oraz funkcje geometryczne w klasie usługowej z jawnymi
sprawdzeniami \code{isNil ifTrue:ifFalse:} w Smalltalku. W Prologu emuluje konstrukcję
if-then-else odcięciami, a nawet zgaduje kolejność argumentów za pomocą klauzul chronionych
odcięciem. Mówiąc wprost, model najpierw pisze algorytm, a dopiero potem dopasowuje go do języka.
Sprawdza się to w językach imperatywnych, gdzie algorytm \emph{jest} idiomem, oraz w językach
funkcyjnych, gdzie zwinięcia i dopasowanie wzorców naturalnie odwzorowują te algorytmy. Sprawdza
się gorzej tam, gdzie paradygmat wymaga innej dekompozycji (obiekty, relacje).

\begin{table}[htbp]
  \centering
  \caption{Antywzorce wskazane przez sędziego (po trzy przykłady na język z \file{report.md}).}
  \label{tab:antipatterns}
  \small
  \begin{tabularx}{\textwidth}{@{}lY@{}}
    \toprule
    Język & Antywzorce \\
    \midrule
    C         & ponowna implementacja \code{fabs} zamiast użycia \code{math.h}; jedna długa funkcja łącząca generowanie punktów, usuwanie duplikatów, budowę otoczki i kopiowanie; ręczna pętla kopiująca zamiast \code{memcpy} \\
    Pascal    & własna implementacja quicksorta zamiast sortowania z RTL lub generycznego; zbędne ręczne czyszczenie pól zarządzanych przed \code{Dispose}; alokacja natychmiast zastępowana inną \\
    Java      & statyczne metody pomocnicze zamiast obiektów dziedzinowych (brak typów otoczki/odcinka); zduplikowane prymitywy geometryczne; surowy bufor \code{Point[]} w stylu C z ręczną żonglerką indeksami \\
    Smalltalk & zduplikowana budowa dolnej i górnej otoczki; funkcje geometryczne w klasie usługowej zamiast w obiektach dziedzinowych; \code{isNil ifTrue:ifFalse:} zamiast \code{ifNil:} \\
    Haskell   & częściowa funkcja \code{head} w \code{map head . groupBy}; \code{length xs < 3} zamiast dopasowania wzorca; leniwe \code{foldl} zamiast \code{foldl'} \\
    Scala     & krawędzie budowane arytmetyką indeksów i modulo zamiast złączenia (\code{zip}) z obróconym wektorem; niespójne użycie epsilona; modyfikowalne \code{var}/\code{while} w wyszukiwaniu binarnym zamiast rekursji ogonowej \\
    Prolog    & odcięcie po warunku dla emulacji if-else; \code{append(A0,[_],A)} do usunięcia ostatniego elementu; pobłażliwe zgadywanie kolejności argumentów w \code{create_polygon/3} z użyciem odcięcia \\
    \bottomrule
  \end{tabularx}
\end{table}

\paragraph{Poprawność i idiomatyczność są w dużej mierze niezależne} Java rozwiązała wszystkie
dziesięć zadań, ale za idiomy paradygmatu uzyskała tylko 3,50. C i Scala są w pełni poprawne
i oceniane wysoko, natomiast Smalltalk i Prolog wypadają nisko w obu wymiarach. Rozwiązania, które
przeszły za pierwszym podejściem, uzyskują średnio 3,64, a te, które wymagały naprawy lub pozostały
nierozwiązane, 3,55, co jest różnicą pomijalną. Poprawność i idiomatyczność trzeba zatem mierzyć
osobno.

\paragraph{Idiomatyczność języka podąża za jego popularnością} W kryterium języka prowadzą szeroko
stosowane Scala (4,00), Java (3,80) i Haskell (3,70), a Smalltalk (3,30), Pascal (3,40) i Prolog
(3,40) zamykają stawkę. Jest to spójne z mniejszą ilością danych treningowych dla tych
języków~\cite{cassano2024knowledge}: model słabiej zna ich biblioteki standardowe i konwencje
(\code{ifNil:}, sortowanie z RTL, predykaty biblioteczne).

\paragraph{Wpływ zadania} Ściśle wyspecyfikowane zadania 7--10 uzyskują niższe oceny (3,39--3,50)
niż zadania 1--6 (3,61--3,93; \cref{tab:judge-heatmap} w \cref{app:results}). Narzucają one
dokładne procedury (operacje na B-drzewie według CLRS, reguły remisów w algorytmie Prima, stabilny
podział trójdzielny), które pozostawiają mało miejsca na idiomatyczną przebudowę. Należy o tym
pamiętać przy porównywaniu lub dodawaniu zadań.

\paragraph{Rozdzielczość ocen} Sędzia korzystał z wąskiej części skali. Wszystkie 280 ocen mieści
się między 3 a 5. Ocenę 5 wystawiono tylko siedem razy, zawsze za idiomatyczność paradygmatu (cztery
w Scali, dwie w Haskellu, jedna w C), a żadne kryterium nigdy nie otrzymało 1 ani 2.
Idiomatyczność języka, czytelność i projekt zawsze otrzymywały 3 lub 4. Przy dziesięciu
rozwiązaniach na język zmiana jednej oceny przesuwa średnią kryterium dla języka o 0,1, więc różnic
rzędu 0,1--0,2 nie należy nadinterpretowywać.

\paragraph{Wiarygodność sędziego} W ramach kalibracji oceniono dwukrotnie pięć rozwiązań wzorcowych
i trzy rozwiązania z przebiegu. Wszystkie 32 pary powtórzonych ocen były identyczne. Między
kalibracją, która używała wcześniejszego schematu bez uzasadnień dla poszczególnych kryteriów,
a ostatecznym sędzią z uzasadnieniem przed oceną trzy rozwiązania z przebiegu różniły się w 1 z 12
ocen (czytelność rozwiązania zadania 3 w C, 5 wobec 4). Sędzia jest zatem stabilny. Czy jest
\emph{trafny}, to osobna kwestia. Rozwiązania wzorcowe napisane na potrzeby walidacji zestawów
testów uzyskały średnio 3,05, wyraźnie mniej niż własne rozwiązania modelu \haiku{} dla tych samych
pięciu zadań (3,90). Może to po prostu oznaczać, że wzorce pisano w celu walidacji testów, a nie
jako rozwiązania modelowe. Może to być też preferowanie własnych odpowiedzi: generator i sędzia
pochodzą z tej samej rodziny modeli, a sędziowie LLM, jak wiadomo, faworyzują tekst podobny do
własnego~\cite{panickssery2024selfpreference}. Bez ocen ludzkich nie potrafimy rozstrzygnąć między
tymi wyjaśnieniami (\cref{sec:threats}).

% -----------------------------------------------------------------------------
\section{RQ4: Koszt, wydajność i wykonywanie instrukcji}
\label{sec:cost}

\Cref{tab:efficiency} i \cref{fig:size} podsumowują koszt przebiegu, rozmiar wygenerowanego kodu
i zgodność modelu z wymaganym formatem odpowiedzi.

\begin{table}[htbp]
  \centering
  \caption{Wydajność i wykonywanie instrukcji według języka. LOC = mediana liczby linii kodu
    pierwszej odpowiedzi; tokeny i czas generowania to średnie na wywołanie API (tokeny wyjściowe
    obejmują rozumowanie); czas ewaluacji to średni czas jednego przebiegu testów w Dockerze;
    naruszenia formatu to odpowiedzi ze znacznikami bloków Markdown.}
  \label{tab:efficiency}
  \small
  \begin{tabular}{@{}lrrrrrrr@{}}
    \toprule
    Język & LOC & \makecell[r]{Tokeny\\wyjściowe} & \makecell[r]{Czas\\gen.\ (s)} & \makecell[r]{Czas\\ewal.\ (s)} & Wywołania & \makecell[r]{Koszt\\(USD)} & \makecell[r]{Naruszenia\\formatu} \\
    \midrule
    C         & 248,5 & 6\,014 & 24,1 &  2,2 & 10 & 0,032 &  5/10 \\
    Pascal    & 273,5 & 9\,320 & 35,5 &  1,8 & 13 & 0,069 &  0/13 \\
    Java      & 184,5 & 3\,791 & 14,0 &  2,6 & 16 & 0,037 &  9/16 \\
    Smalltalk & 132,0 & 7\,166 & 29,8 &  1,2 & 15 & 0,062 &  6/15 \\
    Haskell   & 127,0 & 4\,762 & 19,2 & 32,1 & 15 & 0,042 &  1/15 \\
    Scala     & 104,5 & 4\,459 & 17,3 & 29,3 & 12 & 0,030 &  6/12 \\
    Prolog    & 160,0 & 9\,463 & 39,2 & 26,2 & 17 & 0,095 & 10/17 \\
    \midrule
    Łącznie   & --    & 6\,482 & 25,9 & --   & 98 & 0,368 & 37/98 \\
    \bottomrule
  \end{tabular}
\end{table}

\begin{figure}[htbp]
  \centering
  % Two small multiples: code size and output tokens per call, per language.
% Data: results/haiku-v1/attempts.jsonl (code_lines of attempt 0; output_tokens of all calls).
\begin{subfigure}[t]{0.49\textwidth}
\centering
\begin{tikzpicture}
\begin{axis}[
  reportchart, xbar, width=0.93\linewidth, height=5.6cm, bar width=8pt,
  xmin=0, xmax=300, xtick={0,100,200,300},
  xlabel={Mediana liczby linii kodu (próba 0)},
  ytick={1,...,7}, ymin=0.4, ymax=7.6,
  yticklabels={Prolog, Scala, Haskell, Smalltalk, Java, Pascal, C},
  xmajorgrids=true, axis x line*=bottom, axis y line*=left,
  y axis line style={draw=none}, ytick style={draw=none},
]
  \addplot[fill=cat1, draw=none] coordinates {(248.5,7) (273.5,6) (184.5,5) (132,4) (127,3) (104.5,2) (160,1)};
\end{axis}
\end{tikzpicture}
\caption{Rozmiar rozwiązania.}
\label{fig:size-loc}
\end{subfigure}\hfill
\begin{subfigure}[t]{0.49\textwidth}
\centering
\begin{tikzpicture}
\begin{axis}[
  reportchart, xbar, width=0.93\linewidth, height=5.6cm, bar width=8pt,
  xmin=0, xmax=10, xtick={0,2,...,10},
  xlabel={Tokeny wyjściowe na wywołanie, tys.\ (z rozumowaniem)},
  ytick={1,...,7}, ymin=0.4, ymax=7.6,
  yticklabels={Prolog, Scala, Haskell, Smalltalk, Java, Pascal, C},
  xmajorgrids=true, axis x line*=bottom, axis y line*=left,
  y axis line style={draw=none}, ytick style={draw=none},
]
  \addplot[fill=cat1, draw=none] coordinates {(6.014,7) (9.320,6) (3.791,5) (7.166,4) (4.762,3) (4.459,2) (9.463,1)};
\end{axis}
\end{tikzpicture}
\caption{Nakład na generowanie.}
\label{fig:size-tokens}
\end{subfigure}

  \caption{Rozmiar wygenerowanych rozwiązań i nakład na generowanie według języka. Są to dwie
    odrębne miary, stąd dwie odrębne skale.}
  \label{fig:size}
\end{figure}

\paragraph{Zwięzłość} Rozwiązania funkcyjne są najkrótsze. Mediana rozwiązania w Scali ma 105 linii,
a w Haskellu 127, wobec 249 w C i 274 w Pascalu, tj.\ rozwiązania imperatywne są 2--2,6 razy
dłuższe przy tej samej funkcjonalności. Smalltalk (132) i Prolog (160) plasują się pośrodku.

\paragraph{Nakład na generowanie} Liczba tokenów wyjściowych, łącznie z rozumowaniem, jest
najwyższa dla Prologu (9\,463 na wywołanie) i Pascala (9\,320), a najniższa dla Javy (3\,791). Kod
w Pascalu jest długi, ale w Prologu już nie, więc model poświęca więcej rozumowania mniej znanym
językom. Prolog jest najdroższym językiem: 26\,\% kosztu generowania przy jednej siódmej zadań.
Generowanie zajęło łącznie 42~minuty czasu pracy modelu (25,9\,s na wywołanie), a ewaluacja
w Dockerze dodaje 26--32\,s na przebieg dla Haskella, Scali i Prologu z powodu kompilacji
i uruchamiania.

\paragraph{Koszt} 98 wywołań generujących zużyło 501\,231 tokenów wejściowych i 635\,269 wyjściowych
i kosztowało \usd{0,368}. 70 wywołań sędziego zużyło 403\,703 tokeny wejściowe i 72\,482 wyjściowe
i kosztowało \usd{1,532}, około czterokrotnie więcej niż generowanie, ponieważ sędzia jest większym
modelem i czyta specyfikację, kontrakt oraz rubrykę. Cały przebieg kosztował około \usd{1,90}.
Powtórzone próbki są więc w zasięgu: dziesięć próbek na zadanie kosztowałoby około \usd{19}.

\paragraph{Wykonywanie instrukcji} Choć każdy prompt wymaga surowego kodu, 37 z 98 odpowiedzi
(37,8\,\%) otaczało kod znacznikami bloków Markdown: 28 z 70 pierwszych odpowiedzi (40\,\%) i 9 z 28
odpowiedzi naprawczych (32\,\%). Częstość ta silnie zależy od języka, od 0/13 dla Pascala i 1/15 dla
Haskella do 10/17 dla Prologu i 9/16 dla Javy. Nie wynika ze sformułowań szablonów. Szablony dla
Pascala i Prologu jako jedyne nie wymieniają wprost zakazanego znacznika (\code{no ```c blocks}),
a mimo to jeden z nich ma najniższą częstość, a drugi najwyższą. Znaczniki są nieszkodliwe, bo system
wycina kod. Pięć odpowiedzi z prozą bez znaczników (\cref{sec:failure-analysis}) już takie nie jest,
a metryka oparta na znacznikach ich nie liczy. Zgodność z formatem odpowiedzi najlepiej więc mierzyć
na dwa sposoby: przez znaczniki bloków i przez błędy kompilacji spowodowane tekstem niebędącym
kodem.

% =============================================================================
\chapter{Dyskusja}
\label{ch:discussion}
% =============================================================================

\section{Odpowiedzi na pytania badawcze}
\label{sec:answers}

\paragraph{RQ1: poprawność} Paradygmat ma co najwyżej umiarkowany wpływ na poprawność. Model osiągnął
łącznie 80\,\% poprawności za pierwszym podejściem, z rozrzutem między językami od 6/10 (Java) do
10/10 (C). Rozrzut wewnątrz paradygmatów jest równie duży jak między nimi i przy jednej próbce na
zadanie żadna różnica nie jest istotna. Najwyraźniejszy sygnał jest negatywny: paradygmat logiczny
(Prolog) ma najniższą poprawność. Nawet ten sygnał osłabiają jednak dwa defekty w zestawach testów
dla Prologu.

\paragraph{RQ2: samonaprawa} Informacja zwrotna z wykonania jest skuteczna, ale ograniczona.
Pozwoliła naprawić 9 z 14 nieudanych zadań (64\,\%), wszystkie w ciągu dwóch rund; trzecia runda
nigdy nie pomogła. Naprawa się udaje, gdy informacja zwrotna wskazuje defekt (błędy kompilatora,
niespełnione asercje z komunikatami), i utyka, gdy tego nie robi (połknięte wyjątki, defekty ukryte
w kodzie testów). Błędy formatu w odpowiedziach naprawczych zmarnowały rundy w trzech zadaniach.

\paragraph{RQ3: idiomatyczność} Tu paradygmat ma największe znaczenie. Kod funkcyjny uzyskał
najwyższe oceny (3,90), a obiektowy (3,45) i logiczny (3,40) najniższe, przy czym różnica jest
największa w samym kryterium paradygmatu (4,25 wobec 3,25). Domyślny styl modelu jest algorytmiczny
i proceduralny. Pasuje on do języków imperatywnych i dobrze przekłada się na zwinięcia
i dopasowanie wzorców, ale nie modeluje dziedziny jako obiektów ani relacji, chyba że zostanie do
tego zmuszony. Idiomatyczność języka natomiast podąża za popularnością języka, a nie za jego
paradygmatem.

\paragraph{RQ4: praktyczność} Hostowany mały model sprawia, że badanie jest tanie (około \usd{1,90}
za pełny przebieg, z czego większość przypada na sędziego) i szybkie (42 minuty czasu pracy
modelu). Główną przeszkodą praktyczną jest format odpowiedzi. Na nim poległy lokalne modele
o otwartych wagach, a nawet model hostowany dodał znaczniki bloków do 38\,\% odpowiedzi i prozę do
pięciu z nich, co spowodowało połowę wszystkich błędów kompilacji.

\section{Paradygmat czy język?}
\label{sec:paradigm-or-language}

Projekt z dwoma językami na paradygmat miał rozdzielić te dwa efekty, a wyniki pokazują, dlaczego
było to konieczne. Gdybyśmy dla paradygmatu obiektowego wybrali tylko Javę, wyglądałby on na
najsłabszy pod względem poprawności za pierwszym podejściem. Przy samym Smalltalku wyglądałby
przeciętnie. Haskell i Scala różnią się równie mocno. Przy tej skali każde twierdzenie
o poprawności dla paradygmatu jest więc w istocie twierdzeniem o wybranych językach. Inaczej jest
z idiomatycznością: Java i Smalltalk \emph{obie} uzyskują najniższe oceny za idiomy paradygmatu
(3,50 i 3,00), a Haskell i Scala \emph{oba} najwyższe (4,20 i 4,30). Tutaj efekt paradygmatu jest
spójny dla obu języków z każdej pary, co czyni go bardziej wiarygodnym ustaleniem.

Popularność języka, a więc ilość jego danych treningowych, jest czynnikiem zakłócającym, który ten
projekt ogranicza, ale nie eliminuje. Pascal, Smalltalk i Prolog to języki o małych zasobach.
Pascal i Smalltalk nie są systematycznie gorsze pod względem poprawności, ale wszystkie trzy zajmują
ostatnie miejsca w idiomatyczności języka, a Prolog i Pascal wymagały najwięcej tokenów
rozumowania. Jest to zgodne z poglądem, że model łatwiej przenosi wiedzę algorytmiczną między
językami niż konwencje konkretnego języka~\cite{cassano2024knowledge}.

\section{Dlaczego wszędzie kod proceduralny?}
\label{sec:procedural-bias}

Do proceduralnego stylu domyślnego modelu przyczyniają się prawdopodobnie trzy czynniki.
Po pierwsze, zadania są \emph{algorytmami}, a ich podręcznikowe opisy~\cite{cormen2009clrs} to
imperatywny pseudokod. Im bardziej preskryptywna specyfikacja (zadania 7--10), tym niższe oceny
idiomatyczności. Po drugie, kontrakty projektowano równolegle dla wszystkich języków, wychodząc od C.
Struktury wynikowe, wartości statusu i atomy \code{none} są wierne zadaniu, ale nie zawsze naturalne
dla paradygmatu. Po trzecie, dane treningowe języków obiektowych zawierają mnóstwo kodu w stylu klas
narzędziowych, który model odtwarza. Benchmark, który miałby mierzyć \emph{projekt} obiektowy lub
logiczny, potrzebowałby zadań, których naturalnym rozwiązaniem jest model obiektów lub relacji,
na przykład symulacji, systemów regułowych czy problemów z ograniczeniami z kilkoma oddziałującymi
na siebie bytami, a nie algorytmów z jednym punktem wejścia.

\section{Zagrożenia dla trafności}
\label{sec:threats}

\paragraph{Trafność wewnętrzna}
\begin{itemize}
  \item \emph{Jedna próbka na zadanie.} Modele Claude nie przyjmują ustawienia temperatury, więc
    każda odpowiedź jest losowaniem z rozkładu. Przy jednej próbce zmiana wyniku jednego zadania
    zmienia wynik języka o 10 punktów procentowych. System obsługuje powtórzone próbki
    (\code{SAMPLES}) i są one przystępne cenowo (\cref{sec:cost}).
  \item \emph{Defekty benchmarku.} Dwa zestawy testów dla Prologu zawierają defekty
    (\cref{sec:benchmark-defects}). Podajemy surowe liczby, więc poprawność Prologu jest zaniżona.
    Nie można wykluczyć innych, nieodkrytych defektów, zwłaszcza w ręcznie napisanych zadaniach
    1--5, których nie zweryfikowano względem rozwiązania wzorcowego w każdym języku.
  \item \emph{Nierówne specyfikacje.} Specyfikacje zadań 1--5 dostosowywano ręcznie do każdego
    języka i różnią się one długością i szczegółowością. Na przykład specyfikacja problemu
    plecakowego dla Haskella jest krótsza niż dla C (\cref{app:prompts}). Zadania 6--10
    wygenerowano ze wspólnej specyfikacji i są bardziej jednolite.
  \item \emph{Nierówna informacja zwrotna.} Programy testowe różnią się ilością wypisywanej
    diagnostyki. Program testowy w Smalltalku przechwytuje każdy wyjątek, więc samonaprawa
    w Smalltalku jest trudniejsza z przyczyn niezwiązanych z modelem.
  \item \emph{Testy przechodzące na szkielecie.} 17 testów w C przechodzi na niezaimplementowanym
    szkielecie. Nie zmieniają one żadnego werdyktu rozwiązane/nierozwiązane, ale nieznacznie
    zawyżają odsetki zaliczonych testów dla C.
\end{itemize}

\paragraph{Trafność teoretyczna}
\begin{itemize}
  \item \emph{Sędzia z tej samej rodziny.} Sędzia (\sonnet{}) i generator (\haiku{}) pochodzą
    z tej samej rodziny modeli, a sędziowie LLM mogą preferować odpowiedzi podobne do
    własnych~\cite{zheng2023judging,panickssery2024selfpreference}. Niskie oceny rozwiązań
    wzorcowych (\cref{sec:results-idiomaticity}) są zgodne z takim obciążeniem, ale go nie
    dowodzą.
  \item \emph{Brak ludzkiego punktu odniesienia.} Wykazano, że sędzia jest \emph{stabilny}, ale
    nie że jest \emph{trafny}. Próbkę rozwiązań powinni ocenić ludzcy eksperci w każdym języku.
  \item \emph{Ściśnięta skala.} Sędzia używał tylko ocen 3--5, co ogranicza rozdzielczość. Oceny
    porównawcze (szeregowanie dwóch rozwiązań) lub drobniejsza rubryka mogłyby lepiej rozróżniać
    rozwiązania.
  \item \emph{Format i poprawność są splecione.} Pięć błędów kompilacji spowodowała proza
    w odpowiedzi. Liczą się one jako błędy poprawności, choć są błędami wykonywania instrukcji.
\end{itemize}

\paragraph{Trafność zewnętrzna}
\begin{itemize}
  \item \emph{Jeden generator.} Oceniliśmy jeden mały model na jednym poziomie wysiłku. Większe
    modele, inne rodziny i modele o otwartych wagach mogą wykazywać inne wzorce; pilotaż z modelami
    lokalnymi nie dał użytecznych danych (\cref{sec:local-llms}).
  \item \emph{Dobór zadań.} Dziesięć zadań algorytmicznych to wąski wycinek programowania. Zadania
    skupione na projektowaniu systemów, wejściu/wyjściu, współbieżności czy modelowaniu dziedziny
    mogłyby faworyzować inne paradygmaty.
  \item \emph{Wersje języków.} Narzędzia są ustalone (np.\ Scala 2.11, GNU Smalltalk 3.2.5,
    \cref{tab:languages}), a idiomy nowszych wersji wykraczają poza zakres badania.
\end{itemize}

\paragraph{Trafność wnioskowania statystycznego} Podajemy statystyki opisowe i testy dokładne tylko
dla największych różnic. Przy $n = 10$ zadaniach na język i jednej próbce żadne porównanie nie
osiąga istotności na poziomie $\alpha = 0,05$. Właściwą analizą dla większego przebiegu byłyby
powtórzone próbki (pass@$k$) i model efektów mieszanych z zadaniem i językiem jako czynnikami.

\section{Wnioski dla twórców benchmarków międzyparadygmatowych}
\label{sec:lessons}

\begin{enumerate}
  \item \textbf{Unikać nasyconych problemów.} Problemy, które rozwiązuje każdy model, nie mogą
    ujawnić różnic (\cref{sec:humaneval}). Zadania z wieloma funkcjami, ścisłymi specyfikacjami
    i testami obciążeniowymi to umożliwiły.
  \item \textbf{Każdy kontrakt powinien sam się dokumentować.} Nazwy, kolejność i tryby argumentów
    muszą być widoczne w kontrakcie każdego języka, a nie tylko w nagłówku C
    (\cref{sec:benchmark-defects}).
  \item \textbf{Generować zestawy testów z jednej specyfikacji.} Ręczne przenoszenie testów między
    siedmioma językami wprowadza rozbieżności, takie jak błąd o jeden w Prologu. Umiejętność
    oparta na specyfikacji, użyta dla zadań 6--10, pozwoliła tego uniknąć.
  \item \textbf{Błędy powinny być informatywne.} Programy testowe powinny wypisywać wartości
    oczekiwane i rzeczywiste i nie mogą połykać wyjątków, ponieważ jakość informacji zwrotnej
    ogranicza samonaprawę.
  \item \textbf{Mierzyć format odpowiedzi osobno.} Należy liczyć znaczniki bloków i tekst
    niebędący kodem i oddzielać je od poprawności; inaczej porównanie mierzy wykonywanie
    instrukcji.
  \item \textbf{Walidować przed wydaniem pieniędzy.} Tryb bazowy (szkielety nie przechodzą,
    wzorce przechodzą, testy przechodzące na szkielecie są rejestrowane) nic nie kosztuje
    i wychwytuje zepsutą infrastrukturę testową przed pierwszym płatnym wywołaniem. Nie wychwyci
    jednak testu, który jest błędny w ten sam sposób dla szkieletu i dla każdego rozwiązania,
    więc rozwiązanie wzorcowe dla każdego języka pozostaje mocniejszą kontrolą.
\end{enumerate}

% =============================================================================
\chapter{Podsumowanie i dalsze prace}
\label{ch:conclusions}
% =============================================================================

Zapytaliśmy, czy paradygmat programowania wpływa na kod pisany przez LLM. Aby odpowiedzieć na to
pytanie, zbudowaliśmy benchmark dziesięciu zadań algorytmicznych z identycznymi ukrytymi testami
w siedmiu językach z czterech paradygmatów, wraz z systemem, który automatycznie generuje, testuje,
naprawia i ocenia rozwiązania. Pełny przebieg modelu \haiku{} daje następujący obraz:

\begin{itemize}
  \item \textbf{Poprawność zależy bardziej od języka niż od paradygmatu.} 80\,\% zadań rozwiązano za
    pierwszym podejściem, a 93\,\% po samonaprawie. Różnice między dwoma językami jednego
    paradygmatu są równie duże jak różnice między paradygmatami i żadna nie jest istotna przy
    jednej próbce na zadanie. Prolog wypada najsłabiej, częściowo z powodu dwóch defektów w jego
    zestawach testów.
  \item \textbf{Samonaprawa jest skuteczna, ale ograniczona przez informację zwrotną.} Dwie rundy
    naprawy usunęły 64\,\% błędów, a trzecia runda nic nie dodała. Naprawa utykała tam, gdzie
    informacja zwrotna nie ujawniała przyczyny.
  \item \textbf{Idiomatyczność zależy od paradygmatu.} Kod funkcyjny był najbardziej idiomatyczny,
    a obiektowy i logiczny najmniej. Oba języki każdego paradygmatu są tu zgodne, więc jest to
    najbardziej spójny efekt paradygmatu, jaki zaobserwowaliśmy. Domyślny styl modelu jest
    proceduralny.
  \item \textbf{Format odpowiedzi to sprawa pierwszorzędna.} Przez niego lokalne modele o otwartych
    wagach okazały się dla nas bezużyteczne, a nawet model hostowany wstawił prozę do pięciu
    odpowiedzi, co spowodowało połowę wszystkich błędów kompilacji.
  \item \textbf{Zadania w stylu HumanEval są zbyt łatwe}, by porównywać paradygmaty; potrzebne są
    trudniejsze, precyzyjnie wyspecyfikowane zadania.
\end{itemize}

\section{Dalsze prace}

\begin{enumerate}
  \item \textbf{Poprawienie benchmarku.} Poprawić predykat pomocniczy budujący łańcuch w Prologu
    (zadanie 2), udokumentować nazwy i tryby argumentów w kontrakcie zadania 1 dla Prologu, sprawić,
    by program testowy w Smalltalku raportował komunikaty wyjątków, i ponownie uruchomić dotknięte
    zadania.
  \item \textbf{Powtórzone próbki.} Uruchomić od pięciu do dziesięciu próbek na zadanie, aby
    oszacować pass@$k$ i przedziały ufności, a wyniki analizować z zadaniem i językiem jako efektami
    losowymi i stałymi.
  \item \textbf{Więcej modeli.} Ocenić większe modele (Sonnet, Opus) i inne rodziny modeli, w tym
    modele o otwartych wagach udostępniane przez hostowanego dostawcę inferencji, co pozwala uniknąć
    problemów z wydajnością wdrożenia lokalnego.
  \item \textbf{Walidacja sędziego.} Zebrać oceny ludzkich ekspertów dla warstwowej próbki
    rozwiązań, użyć sędziego z innej rodziny modeli i porównać oceny bezwzględne z preferencjami
    parami.
  \item \textbf{Zadania natywne dla paradygmatów.} Dodać zadania, których naturalnym rozwiązaniem
    jest model obiektowy lub zbiór relacji (symulacje, silniki reguł, problemy z ograniczeniami),
    a także kolejne paradygmaty lub języki (np.\ OCaml, Erlang/Elixir, Rust).
  \item \textbf{Ablacje promptów.} Zmierzyć, w jakim stopniu krótkie wskazówki dotyczące paradygmatu
    w promptach (\cref{tab:guidance}) zmieniają idiomatyczność, oraz czy ściślejszy kontrakt
    odpowiedzi (np.\ pole z kodem w ustrukturyzowanej odpowiedzi) eliminuje błędy formatu.
\end{enumerate}


% ------------------------------------------------------------------ część końcowa
\clearpage
\phantomsection
\addcontentsline{toc}{chapter}{Bibliografia}
\bibliographystyle{plainnat}
\bibliography{references}

\appendix
% =============================================================================
\chapter{Prompty}
\label{app:prompts}
% =============================================================================

Ten dodatek przytacza szablony promptów wszystkich języków (\file{<Language>/README.md}), jeden
kompletny wygenerowany prompt początkowy oraz rzeczywisty prompt naprawczy z przebiegu. Prompty
podajemy w oryginalnym brzmieniu angielskim, w jakim trafiały do modelu. \Cref{sec:prompts} opisuje,
jak wypełniane są szablony.

\section{Szablony promptów początkowych}
\label{app:templates}

Wszystkie szablony mają strukturę z \cref{lst:prompt-initial}: rola i cel, symbol zastępczy
\code{[TASK-SPECIFIC REQUIREMENTS]}, wymagania ogólne języka i symbol zastępczy kontraktu. Szablon
dla C pokazuje \cref{lst:prompt-initial}; pozostałe sześć przytaczamy poniżej.

\begin{lstlisting}[style=prompt,caption={Szablon promptu początkowego: Pascal.}]
You are a Pascal programming expert. Your task is to implement the required functions in accordance with the provided unit interface.

[TASK-SPECIFIC REQUIREMENTS]

General Requirements:
- Use idiomatic procedural Pascal.
- Manage dynamically allocated structures safely.
- Return ONLY raw Pascal code (no Markdown tags, no side text).
- Do not provide test runner code.

[INSERT UNIT INTERFACE HERE]
\end{lstlisting}

\begin{lstlisting}[style=prompt,caption={Szablon promptu początkowego: Java.}]
You are a Java programming expert. Your task is to implement the required methods in accordance with the provided class or interface.

[TASK-SPECIFIC REQUIREMENTS]

General Requirements:
- Use idiomatic Java.
- Use object-oriented design and model the domain with classes.
- Prefer standard library collections and utilities where appropriate.
- Return ONLY raw Java code (no Markdown tags, no ```java blocks, no side text). Your code will be directly saved to a .java file and compiled. Do not provide a main() method or test code.

[INSERT CLASS OR INTERFACE CONTENT HERE]
\end{lstlisting}

\begin{lstlisting}[style=prompt,caption={Szablon promptu początkowego: Smalltalk.}]
You are a GNU Smalltalk programming expert. Your task is to implement the required methods in accordance with the provided class hierarchy.

[TASK-SPECIFIC REQUIREMENTS]

General Requirements:
- Use idiomatic GNU Smalltalk (gst) syntax.
- Use object-oriented design and model the domain with classes and message passing.
- Do not redefine built-in classes (e.g. Point, Set, Interval); use them as provided.
- Prefer standard Smalltalk collections and idioms where appropriate.
- Return ONLY raw Smalltalk code (no Markdown tags, no ```smalltalk blocks, no side text). Your code will be directly saved to a .st file and loaded by gst. Do not provide test code.

[INSERT INTERFACE FILE CONTENT HERE]
\end{lstlisting}

\begin{lstlisting}[style=prompt,caption={Szablon promptu początkowego: Haskell.}]
You are a Haskell programming expert. Your task is to implement the required functions in accordance with the provided module interface.

[TASK-SPECIFIC REQUIREMENTS]

General Requirements:
- Use idiomatic, pure Haskell.
- Prefer immutable data, algebraic data types, and total functions where practical.
- Model absence with `Maybe` instead of null pointers.
- Return ONLY raw Haskell code (no Markdown tags, no ```haskell blocks, no side text). Your code will be directly saved to a .hs file and compiled.
- Do not provide test runner code.

[INSERT MODULE INTERFACE HERE]
\end{lstlisting}

\begin{lstlisting}[style=prompt,caption={Szablon promptu początkowego: Scala.}]
You are a Scala programming expert. Your task is to implement the required functions in accordance with the provided module interface.

[TASK-SPECIFIC REQUIREMENTS]

General Requirements:
- Use idiomatic Scala.
- Prefer immutable data structures and pure functions where practical.
- Model absence with `Option` instead of null values.
- Return ONLY raw Scala code (no Markdown tags, no ```scala blocks, no side text). Your code will be directly saved to a `.scala` file and compiled.
- Do not provide test runner code.

[INSERT MODULE INTERFACE HERE]
\end{lstlisting}

\begin{lstlisting}[style=prompt,caption={Szablon promptu początkowego: Prolog.}]
You are a Prolog programming expert. Your task is to implement the required predicates in accordance with the provided module interface.

[TASK-SPECIFIC REQUIREMENTS]

General Requirements:
- Use idiomatic SWI-Prolog.
- Prefer pure predicates and relational style where practical.
- Model failure through predicate failure rather than exceptions.
- Return ONLY raw Prolog code (no Markdown tags, no side text). Your code will be directly saved to a .pl file and executed.
- Do not provide test runner code.

[INSERT MODULE INTERFACE HERE]
\end{lstlisting}

\section{Szablony promptów naprawczych}

Szablony naprawcze wszystkich języków mają postać z \cref{lst:prompt-followup}. Różnią się tylko
etykietą symbolu zastępczego logu i dwoma pierwszymi wymaganiami dotyczącymi poprawki
(\cref{tab:followup}). Szablony imperatywne i obiektowe akcentują bezpieczeństwo pamięci i zasobów,
a funkcyjne i logiczne powtarzają prośbę o idiomatyczny kod. Tę asymetrię należy mieć na uwadze
przy porównywaniu idiomatyczności naprawionych rozwiązań.

\begin{table}[htbp]
  \centering
  \caption{Części szablonów naprawczych specyficzne dla języka.}
  \label{tab:followup}
  \small
  \begin{tabularx}{\textwidth}{@{}llYY@{}}
    \toprule
    Język & Symbol zastępczy logu & Wymaganie 1 & Wymaganie 2 \\
    \midrule
    C         & \code{[ERROR LOGS / VALGRIND OUTPUT]} & usunąć główny problem logiczny lub pamięciowy & zapewnić bezpieczne zarządzanie pamięcią \\
    Pascal    & \code{[ERROR LOGS]}                   & usunąć główny problem logiczny lub pamięciowy & zapewnić bezpieczne zarządzanie pamięcią \\
    Java      & \code{[ERROR LOGS / COMPILER OUTPUT]} & usunąć główny problem logiczny lub problem zarządzania zasobami & zapewnić bezpieczną obsługę zasobów \\
    Smalltalk & \code{[ERROR LOGS / GST OUTPUT]}      & usunąć główny problem logiczny & zapewnić bezpieczną obsługę argumentów nil i przypadków brzegowych \\
    Haskell   & \code{[ERROR LOGS]}                   & usunąć główny problem logiczny & zachować idiomatyczność rozwiązania w Haskellu \\
    Scala     & \code{[ERROR LOGS]}                   & usunąć główny problem logiczny & zachować idiomatyczność rozwiązania w Scali \\
    Prolog    & \code{[ERROR LOGS]}                   & usunąć główny problem logiczny & zachować idiomatyczność rozwiązania w Prologu \\
    \bottomrule
  \end{tabularx}
\end{table}

\section{Kompletny wygenerowany prompt}

\Cref{lst:rendered} przedstawia dokładny prompt początkowy, który \haiku{} otrzymał dla zadania 4
(plecak) w Haskellu: szablon z wypełnionymi wymaganiami specyficznymi dla zadania i modułem-szkieletem.
Można go odtworzyć poleceniem \code{python harness/run.py --dry-run}.

\begin{lstlisting}[style=prompt,caption={Wygenerowany prompt początkowy, Haskell, zadanie 4.},label={lst:rendered}]
You are a Haskell programming expert. Your task is to implement the required functions in accordance with the provided module interface.

- Implement `maxValue weights values capacity`: Return the maximum value attainable with sum of selected weights ≤ capacity. Use DP. Return 0 if null weights or capacity <= 0.
- Implement `solve weights values capacity`: Return Result with resultMaxValue, resultSelected (length n, entries 0/1), resultTotalWeight.
- DO NOT use greedy heuristics — they fail on w=[10,20,30], v=[60,100,120], cap=50 (optimum=220, greedy=160).
- Selection invariants: sum of selected weights ≤ capacity; sum of selected values = maxValue.
- Use idiomatic Haskell (immutable lists / arrays).

General Requirements:
- Use idiomatic, pure Haskell.
- Prefer immutable data, algebraic data types, and total functions where practical.
- Model absence with `Maybe` instead of null pointers.
- Return ONLY raw Haskell code (no Markdown tags, no ```haskell blocks, no side text). Your code will be directly saved to a .hs file and compiled.
- Do not provide test runner code.

module Knapsack
  ( Result (..)
  , maxValue
  , solve
  ) where

data Result = Result
  { resultMaxValue    :: Int
  , resultSelected    :: [Int]
  , resultTotalWeight :: Int
  } deriving (Eq, Show)

maxValue :: [Int] -> [Int] -> Int -> Int
maxValue _ _ _ = error "Knapsack.maxValue not implemented"

solve :: [Int] -> [Int] -> Int -> Result
solve _ _ _ = error "Knapsack.solve not implemented"
\end{lstlisting}

\section{Rzeczywisty prompt naprawczy}

Pierwsza odpowiedź na ten prompt kończyła się zabłąkanym znacznikiem \code{</br>}
(\cref{tab:compile-errors}). \Cref{lst:followup-real} to prompt naprawczy, który odesłał system.
Fragment logu to wyjście kompilatora, a ponieważ żaden test się nie uruchomił, linia z niezaliczonymi
testami to odnotowuje. Następna odpowiedź przeszła wszystkie 20 testów.

\begin{lstlisting}[style=prompt,caption={Prompt naprawczy wysłany po błędzie kompilacji w zadaniu 4 w Haskellu (próba 1).},label={lst:followup-real}]
Your previous Haskell implementation failed the evaluation. Your task is to analyze the error and provide a corrected version of the code.

ghc -Wall -Wextra -g -itest-4 -o test_runner test-4/Tests.hs
[1 of 2] Compiling Knapsack         ( test-4/Knapsack.hs, test-4/Knapsack.o )

test-4/Knapsack.hs:66:1: error: parse error on input `</'
   |
66 | </br>
   | ^^
make: *** [Makefile:7: test] Error 1

Failing tests: all (no test results were produced, see above).

Requirements for the fix:
1. Identify and resolve the core logical issue causing the failure.
2. Keep the solution idiomatic for Haskell.
3. Return ONLY the complete, raw, fixed Haskell code for `Knapsack.hs`.
4. Do NOT include Markdown formatting, do NOT include ```haskell tags, and do NOT write any explanations. Just the code.
\end{lstlisting}

% =============================================================================
\chapter{Sędzia LLM: prompt, rubryka i przykład}
\label{app:judge}
% =============================================================================

\section{Prompt sędziego}

Sędzia otrzymuje prompt systemowy (\cref{lst:judge-system}), po którym następuje rubryka, oraz
wiadomość użytkownika (\cref{lst:judge-user}) ze specyfikacją zadania, kontraktem i kodem. Symbole
zastępcze w nawiasach klamrowych są wypełniane dla każdego rozwiązania. Sędzia nigdy nie otrzymuje
wyników testów, liczby prób ani nazwy modelu generującego. Prompty i rubrykę podajemy
w oryginalnym brzmieniu angielskim; \cref{app:rubric} zawiera ich polskie omówienie.

\begin{lstlisting}[style=prompt,caption={Prompt systemowy sędziego (\file{harness/judge.py}).},label={lst:judge-system}]
You are an expert reviewer of {language} code and of the {paradigm} programming paradigm. You assess how idiomatic a solution to a programming task is. Apply the rubric below and answer only with the requested JSON.

Before the scores, write a rationale for each criterion: point to the concrete code that decided the score (function, type or predicate names, specific constructs) and say what would have earned one point more, or why the code already deserves a 5. A reader must be able to check every score against the code from its rationale alone.

{general rubric sections + the section of this language}
\end{lstlisting}

\begin{lstlisting}[style=prompt,caption={Struktura wiadomości użytkownika dla sędziego.},label={lst:judge-user}]
Language: {language} ({paradigm} paradigm)

## Task specification given to the author
{task-specific requirements from test-N.md}

## Contract given to the author ({contract file}, not written by the author)
   -- or, in single-file languages --
## Stub module given to the author ({module file})
The author replaced the stub bodies; type declarations and signatures come from it.
```{language tag}
{contract}
```

## Code to evaluate ({target file})
```{language tag}
{code}
```
\end{lstlisting}

\section{Rubryka}
\label{app:rubric}

Rubryka (\file{harness/rubrics.md}) składa się z sekcji ogólnych, które otrzymuje każda ocena, oraz
z sekcji dla każdego języka, z których dołączana jest tylko sekcja ocenianego języka. Poniżej
podajemy jej treść w tłumaczeniu.

\paragraph{Skala}
\begin{itemize}
  \item[\textbf{5}] Tak napisałby to doświadczony praktyk tego języka i paradygmatu; nie ma czego
    zmieniać.
  \item[\textbf{4}] Ogólnie idiomatyczny, z kilkoma drobnymi odstępstwami od powszechnej praktyki.
  \item[\textbf{3}] Miejscami zgodny ze stylem języka, ale z wyraźnymi nieidiomatycznymi
    fragmentami.
  \item[\textbf{2}] Napisany przeważnie tak, jakby był w innym języku lub paradygmacie; idiomy tego
    języka są w dużej mierze ignorowane.
  \item[\textbf{1}] Walczy z językiem: wszędzie obce mu konstrukcje lub niewłaściwe użycie jego
    podstawowych cech.
\end{itemize}

\paragraph{Kryteria}
\begin{itemize}
  \item \code{paradigm_idiomaticity}: jak dobrze kod wyraża rozwiązanie w paradygmacie języka
    (proceduralnym, obiektowym, funkcyjnym, logicznym): dekompozycja, przepływ danych i struktury
    sterujące preferowane przez ten paradygmat.
  \item \code{language_idiomaticity}: wykorzystanie cech danego języka, jego biblioteki
    standardowej i konwencji (nazewnictwo, reprezentacja błędów i braku wartości, kolekcje,
    formatowanie).
  \item \code{readability}: jak łatwo kompetentny czytelnik tego języka śledzi kod: nazwy, rozmiar
    funkcji, struktura, komentarze tam, gdzie pomagają (a nie wszędzie).
  \item \code{design}: podział na funkcje/typy/klasy/predykaty, dobór struktur danych, rozdzielenie
    odpowiedzialności, unikanie powtórzeń; struktura, a nie poprawność.
\end{itemize}

\paragraph{Reguły}
\begin{itemize}
  \item Nie oceniać poprawności: poprawne, ale nieidiomatyczne rozwiązanie powinno dostać niską
    ocenę, a idiomatyczne rozwiązanie z błędem może dostać wysoką.
  \item Oceniać tylko to, co napisał autor; deklaracje i sygnatury skopiowane z kontraktu są dane.
  \item Nie nagradzać długości ani dodatkowych funkcji i nie karać krótkiego rozwiązania za to, że
    jest krótkie.
  \item Wymagania wynikające z zadania (ustalone sygnatury, struktury wynikowe w stylu C, typy
    wyliczeniowe statusów, atomy \code{none}, kolekcje Smalltalka indeksowane od 1) nie są
    antywzorcami.
  \item \code{anti_patterns}: krótkie, konkretne etykiety nieidiomatycznych konstrukcji faktycznie
    obecnych w kodzie (np.\ „pętle po indeksach zamiast zwinięć”, „odcięcie użyte do emulacji
    if-else”); bez drobiazgów stylistycznych.
  \item \code{justification}: od dwóch do czterech zdań wskazujących konkretne fragmenty kodu.
\end{itemize}

\paragraph{Sekcje językowe} \Cref{tab:rubric-langs} streszcza oczekiwania wymienione dla każdego
języka.

\begin{table}[htbp]
  \centering
  \caption{Oczekiwania rubryki specyficzne dla języka (w skrócie).}
  \label{tab:rubric-langs}
  \small
  \begin{tabularx}{\textwidth}{@{}lY@{}}
    \toprule
    Język & Oczekiwania \\
    \midrule
    C         & jasna własność pamięci z odpowiadającymi funkcjami \code{free_*} i brakiem wycieków na ścieżkach błędów; sprawdzane wyniki \code{malloc}/\code{calloc}; funkcje pomocnicze \code{static}, \code{const} dla wskaźników tylko do odczytu; małe funkcje o jednym przeznaczeniu; biblioteka standardowa (\code{qsort}, \code{memcpy}, \code{snprintf}) zamiast własnych implementacji \\
    Pascal    & ustrukturyzowane procedury ze świadomie dobranymi parametrami \code{const}/\code{var}; spójne \code{New}/\code{Dispose} lub \code{SetLength}; narzędzia RTL Free Pascala (\code{Length}, \code{High}, \code{Exit(value)}, \code{Math}, \code{SysUtils}); czytelny przepływ sterowania, bez \code{goto} \\
    Java      & zachowanie w dobrze nazwanych klasach i prywatnych metodach pomocniczych z hermetyzowanym stanem; framework kolekcji i standardowe API; \code{final}, typy generyczne, typy wyliczeniowe, niemodyfikowalność; brak połykanych wyjątków; bez statycznego „śmietnika” narzędziowego i kodu w stylu C \\
    Smalltalk & przekazywanie komunikatów i polimorfizm z małymi metodami o czytelnych intencjach; protokoły kolekcji (\code{do:}, \code{collect:}, \code{select:}, \code{detect:ifNone:}, \code{inject:into:}); idiomatyczne bloki (\code{ifNil:}, \code{and:}); logika dziedzinowa w obiektach; bez redefiniowania klas wbudowanych \\
    Haskell   & czyste, totalne funkcje; \code{Maybe}/\code{Either}/ADT zamiast \code{error} lub funkcji częściowych; dopasowanie wzorców, strażnicy, zwinięcia; \code{Data.Map}/\code{Data.Set}/\code{Data.Array} do wyszukiwania; ścisłość tam, gdzie ma znaczenie (\code{foldl'}); bez emulacji stylu imperatywnego \\
    Scala     & \code{val} i niemodyfikowalne kolekcje, wyrażenia zamiast instrukcji; \code{Option}/\code{Either}/zapieczętowane cechy/klasy przypadków; kombinatory kolekcji i for-comprehensions; \code{@tailrec}; bez nulli w stylu Javy i \code{return}; Scala 2.11 \\
    Prolog    & predykaty relacyjne z rozgałęzianiem w głowach klauzul; akumulatory lub predykaty biblioteczne (\code{foldl/4}, \code{maplist/N}, \code{aggregate_all/3}); tylko „zielone” lub uzasadnione odcięcia; bez \code{assert}/\code{retract} dla zwykłego stanu; porażka lub \code{none} zgodnie ze specyfikacją \\
    \bottomrule
  \end{tabularx}
\end{table}

\section{Przykładowa ocena}

\Cref{lst:judgement} przedstawia kompletną ocenę rozwiązania zadania 1 w C (oceny 4/4/4/4). Każde
uzasadnienie wskazuje konkretne funkcje i konstrukcje oraz to, co dałoby wyższą ocenę, zgodnie
z wymaganiem schematu. Ocenę podajemy w oryginalnym brzmieniu.

\begin{lstlisting}[language=JSON,numbers=none,caption={Ocena rozwiązania zadania 1 w C (\file{results/haiku-v1/judge/C-test-1-s1-a0.json}, bez metadanych).},label={lst:judgement}]
{
  "rationale": {
    "paradigm_idiomaticity": "The code is clean procedural C: small static helpers (geo_cross, geo_on_segment, geo_segments_intersect), explicit loops, data passed through structs, and straightforward control flow with early returns. Monotone-chain hull built on qsort is the standard imperative approach. Cleanup uses a single exit path in places, with early-break loops. A 5 would require fewer inline loops in the big functions, e.g. extracting hull construction and the 'inside' test into helpers.",
    "language_idiomaticity": "Uses static helpers, qsort with a comparator, size_t for sizes, calloc/malloc checks, and a bounded strncpy wrapper that guarantees termination. Naming is consistent with a geo_ prefix. Minor departures: [...] the cast of calloc/malloc results is unnecessary in C; geo_abs reimplements fabs from <math.h>; the `for (size_t i = u - 1; i-- > 0; )` idiom is a bit cryptic; [...]",
    "readability": "Names are clear and comments mark the hull phases and the collision steps. Functions are mostly short, but calculate_clean_minkowski_sum is long, mixing point generation, dedupe, hull and result copy. [...] Splitting it into helpers would earn a 5.",
    "design": "Good separation: geometry primitives, name copy, allocation/free pairs, and collision implemented via the Minkowski sum of obstacle and negated robot, as the spec requires. Ownership is clear and error paths free pts/hull/neg/cspace. Weaknesses: [...] hull construction is not factored into a reusable convex_hull function; [...] the final copy from hull to result could be a memcpy."
  },
  "paradigm_idiomaticity": 4,
  "language_idiomaticity": 4,
  "readability": 4,
  "design": 4,
  "anti_patterns": [
    "reimplementing fabs (geo_abs) instead of using math.h",
    "long function mixing point generation, dedupe, hull building and result copy",
    "manual loop copy instead of memcpy for hull to result vertices",
    "unnecessary casts on malloc/calloc results"
  ],
  "justification": "The code is solid procedural C with static helpers, qsort, checked allocations and consistent cleanup, and the collision check follows the C-space approach via the Minkowski sum with the negated robot. Small deductions come from the long calculate_clean_minkowski_sum, the hand-rolled geo_abs and copy loop, and the lack of a separate convex-hull helper. Overall it is idiomatic and easy to follow."
}
\end{lstlisting}

Dla kontrastu rozwiązanie zadania 4 w Prologu (oceny 4/3/3/3) zostało opisane jako „czyste,
pozbawione \code{assert} programowanie dynamiczne, które przekazuje niemodyfikowalną tablicę
pojemności przez \code{fold_items/step/combine}”, z antywzorcami „ręcznie napisane
\code{fold_items} zamiast \code{foldl/4}”, „łańcuchy if-then-else w \code{combine/4} zamiast
rozgałęziania w głowach klauzul” oraz „przesuwanie list za pomocą \code{append} w celu emulacji
indeksowania tablic”.

% =============================================================================
\chapter{Szczegółowe wyniki}
\label{app:results}
% =============================================================================

Tabele tego dodatku są generowane z \file{results/haiku-v1} przez skrypt
\file{Report/tools/make_tables.py}.

\section{Idiomatyczność według języka i zadania}

\begin{table}[htbp]
  \centering
  \caption{Średnia ocena sędziego (średnia z czterech kryteriów) według języka i zadania.
    Ciemniejsze komórki oznaczają kod bardziej idiomatyczny.}
  \label{tab:judge-heatmap}
  \footnotesize
  \setlength{\tabcolsep}{4pt}
  \begin{tabular}{@{}l*{10}{c}c@{}}
    \toprule
    Język & Z1 & Z2 & Z3 & Z4 & Z5 & Z6 & Z7 & Z8 & Z9 & Z10 & Średnia \\
    \midrule
    \tabinput{tables/judge_heatmap}
    \bottomrule
  \end{tabular}
\end{table}

\begin{table}[htbp]
  \centering
  \caption{Ocena idiomatyczności paradygmatu według języka i zadania. Smalltalk otrzymał 3
    w każdym zadaniu.}
  \label{tab:paradigm-scores}
  \footnotesize
  \setlength{\tabcolsep}{6pt}
  \begin{tabular}{@{}l*{10}{c}@{}}
    \toprule
    Język & Z1 & Z2 & Z3 & Z4 & Z5 & Z6 & Z7 & Z8 & Z9 & Z10 \\
    \midrule
    \tabinput{tables/paradigm_scores}
    \bottomrule
  \end{tabular}
\end{table}

\section{Wszystkie zadania}

\begingroup
\footnotesize
\setlength{\tabcolsep}{3.5pt}
\begin{longtable}{@{}lrlrrcrrccccc@{}}
  \caption{Wyniki wszystkich 70 zadań. \emph{Pierwszy wynik}: kategoria próby 0; \emph{Pierwsza} /
    \emph{Końc.}: liczba zaliczonych testów (z 20); \emph{Rozw.\ w}: próba, w której rozwiązano
    zadanie (0 = pierwsza odpowiedź, \textemdash{} = nierozwiązane); \emph{Pr.}: wykorzystane próby;
    \emph{Tokeny}: tokeny wyjściowe wszystkich prób (łącznie z rozumowaniem); P/J/C/D: oceny
    sędziego za idiomy paradygmatu, idiomy języka, czytelność i projekt.}
  \label{tab:full-results}\\
  \toprule
  Język & Z & Pierwszy wynik & Pierwsza & Końc. & Rozw.\ w & Pr. & Tokeny & P & J & C & D & Średnia \\
  \midrule
  \endfirsthead
  \multicolumn{13}{@{}l}{\small\itshape \Cref{tab:full-results} -- ciąg dalszy}\\
  \toprule
  Język & Z & Pierwszy wynik & Pierwsza & Końc. & Rozw.\ w & Pr. & Tokeny & P & J & C & D & Średnia \\
  \midrule
  \endhead
  \bottomrule
  \endfoot
  \tabinput{tables/full_results}
\end{longtable}
\endgroup


\end{document}
